\documentclass{article}
%\documentclass[10pt,conference]{IEEEtran}
\usepackage{amsmath,amssymb,amsthm, amsfonts, mathrsfs}
\usepackage{hyperref}
\usepackage{xcolor}
\usepackage{textcomp}
\usepackage{enumerate} \usepackage{amsmath}
\usepackage{calligra,indentfirst,epsfig}
\calligra
\addtolength{\hoffset}{-0.55cm}\setlength
{\hoffset}{-0.5in}
\addtolength{\voffset}{-0.45cm}\setlength{\voffset}{-0.5in}
\addtolength{\textwidth}{1.5cm} \addtolength{\textheight}{1.7cm}
\setlength{\textwidth}{6.8in} \setlength{\textheight}{9 in}
\setlength {\topmargin}{0.01in} \setlength{\evensidemargin}{0.45in}
\setlength{\oddsidemargin}{0.45in} \setlength {\columnsep}{6mm}
%\newcommand{\co}{\mathcal{C}}
%\newcommand{\Fpm}{{\mathbb{F}_{p^m}}}
%\newcommand{\Fqn}{{\mathbb{F}_{q}}^{n}}
%\newcommand{\Fq}{\mathbb{F}_{q}}
%\newcommand{\code}{\mathscr{C}}
%\newcommand{\f}{\frac{m}{2}}
%\newcommand{\I}{\mathcal{I}}
%\newcommand{\F}{\mathcal{F}}
%\newcommand{\Sr}{\mathcal{S}}
%\newcommand{\fe}{\frac{m+1}{2}}
%\newcommand{\boldm}[1] {\mathversion{bold}#1\mathversion{normal}}

%\newcommand{\textsubscript}[1] {_\text#1}

\hypersetup{
           breaklinks=true,   % splits links across lines
           colorlinks=true,   % displays links as colored text instead of blocks
           pdfusetitle=true,  % \title and \author values into pdf metadata
                              % etc.
        }

%\hypersetup{
%           breaklinks=true,   % splits links across lines
%           colorlinks=true,   % displays links as colored text instead of blocks
%           pdfusetitle=true,  % \title and \author values into pdf metadata
%                              % etc.
%        } 
   
       

\DeclareMathOperator{\tr}{tr}
\DeclareMathOperator{\Tr}{Tr}
\DeclareMathOperator{\rank}{rank}

\def\F{\mathbb{F}}
\def\Z{\mathbb{Z}}

\newtheorem{thm}{Theorem}[section]
\newtheorem{lem}{Lemma}[section]
\newtheorem{cor}{Corollary}[section]
\newtheorem{prop}{Proposition}[section]
\newtheorem{notation}{Notations}[section]
\theoremstyle{remark}
\newtheorem{rem}{Remark}[section]
\newtheorem{ex}{Example}[section]
\theoremstyle{definition}
\newtheorem{defin}{Definition}[section]
\allowdisplaybreaks
\linespread{1.2}
\numberwithin{equation}{section}
\usepackage{breakurl}
\usepackage{url}
\begin{document} 
 
\title{\bf {Entanglement-Assisted Quantum Error-Correcting Codes over Local Frobenius Rings}}
 
\author{Tania Sidana and Navin Kashyap\thanks{The authors are with the Department of Electrical Communication Engineering, Indian Institute of Science, Bangalore 560012. Email: \{taniasidana,nkashyap\}@iisc.ac.in. This work was supported an IISc-IoE postdoctoral fellowship awarded to the first author.}}

\date{}
 
\maketitle 

 \begin{abstract} In this paper, we provide a framework for constructing entanglement-assisted quantum error-correcting codes (EAQECCs) from classical additive codes over a finite commutative local Frobenius ring. We derive a formula for the minimum number of entanglement qudits required to construct an EAQECC from a linear code over the ring $\mathbb{Z}_{p^s}$. This is used to obtain an exact expression for the minimum number of entanglement qudits required to construct an EAQECC from an additive code over a Galois ring, which significantly extends known results for EAQECCs over finite fields. \end{abstract}

 {\bf Keywords:} { quantum error correction; entanglement-assisted codes; Galois rings; Frobenius rings}

 \section{Introduction}
Quantum error-correcting codes (QECCs) protect quantum states against decoherence caused by the interaction between quantum states and their environment. The stabilizer framework, proposed by Gottesman \cite{gottesman}, is one of the main mechanisms for constructing QECCs. The construction is based on abelian subgroups of the Pauli group, and the resulting QECCs are called quantum stabilizer codes. The stabilizer framework encompasses the first QECC constructed by Shor \cite{shor}, as well as the construction from classical error-correcting codes, discovered independently by Calderbank and Shor \cite{calderbank} and Steane \cite{steane}. The latter construction, now known as the Calderbank-Shor-Steane (CSS) construction, uses dual-containing (or self-orthogonal) classical codes to form QECCs. 

Originally developed for qubits, the stabilizer framework was subsequently extended to higher-dimensional qudit spaces. The (Pauli) error group in this case is generated by unitary operators whose actions on qudits are defined by the algebra of an underlying finite field or ring. The extension of the stabilizer framework to error groups defined via finite fields was executed by Ashikhmin and Knill \cite{ashikhmin}, and Ketkar et al.\ \cite{ketkar}, while the same was done for finite commutative Frobenius rings by Nadella and Klappenecker \cite{nadella}, and Gluesing-Luerssen and Pllaha \cite{luerssen}.

The stabilizer formalism was extended in a different direction by Brun et al.\ \cite{brun}, who gave a method of constructing QECCs (over qubits) from \emph{non-abelian} subgroups of the Pauli group. The idea here was to add more dimensions to the Pauli group, so as to introduce extra degrees of freedom that can be used to ``abelianize'' the original non-abelian subgroup. The code construction required the existence of a small number of pre-shared entanglement qubits between the sender and receiver, where the receiver-end qubits are assumed to be error-free throughout. As a consequence, these codes were called entanglement-assisted quantum error correcting codes (EAQECCs). Within this framework, \emph{any} classical binary linear code can be used as the starting point for constructing an EAQECC. Wilde and Brun \cite{wilde} determined the minimal number of shared qubits required to construct an EAQECC starting from a given binary linear code, and more generally, starting from a given non-abelian subgroup of the Pauli group. 

The theory of EAQECCs extends readily to qudit spaces for which the (Pauli) error groups are defined by finite fields. Indeed, Wilde and Brun observe in \cite[Remark~1]{wilde} that their formula for the minimum number of shared qudits also applies to EAQECCs constructed from linear codes over any prime field; a formal proof of this was given by Luo et al.\ \cite{luo}. Later, Galindo et al.\ \cite{galindo} verified that this formula also applied to EAQECCs obtained from linear codes over an arbitrary finite field $\F_q$. More recently, Nadkarni and Garani \cite{nadkarni} derived an analogous formula for EAQECCs constructed from $\F_p$-additive codes over $\F_q$, where $p$ is the characteristic of $\F_q$.

In this paper, we extend the EAQECC formalism to qudit spaces on which error actions are defined by finite local commutative Frobenius rings. This enables us to construct EAQECCs starting from classical additive codes over such rings, which are overall a much richer class of (classical) codes. It must be pointed out here that Lee and Klappenecker \cite{lee} have previously studied EAQECCs over finite commutative (but not necessarily local) Frobenius rings. However, their EAQECC construction relies crucially on Theorem~5 of their paper, in the proof of which we found a gap that could not readily be filled.\footnote{In the proof of \cite[Theorem 5]{lee}, Lee and Klappenecker replace $\mathbf{w}_k$ with $\mathbf{w}'_{k-2} = e_{k,i} \mathbf{w}_k - \cdots $. However, this may not result in a basis of $R^{2n},$ as $e_{k,i}$ need not be a unit in the ring $R$.}  By restricting our attention further to local rings, we are able to furnish a proof of the result (Theorem~\ref{sym}) needed for the construction of EAQECCs. Thus, one of the contributions of this paper is to provide a coding-theoretic framework to construct EAQECCs over finite commutative local Frobenius rings from first principles. 

We also attempt to obtain the minimum number of pre-shared entanglement qudits required to construct an EAQECC, starting from an additive code over such a ring. We succeed in deriving a formula for this number in the special case of EAQECCs over finite Galois rings. To get to the answer, we had to first derive it for the basic case of the integer rings $\Z_{p^s},$ which itself turned out to be a non-trivial task. 

This paper is organized as follows. In Section \ref{sec2}, we establish the basic definitions and notation needed to describe the construction of quantum stabilizer codes and EAQECCs. This section also contains statements of our main results.  In Section \ref{sec3}, we provide the means to construct EAQECCs from any additive code over a finite commutative local Frobenius ring. In Section \ref{sec4}, we derive an explicit form of the minimum number of pre-shared entanglement qudits required to construct an EAQECC over a ring, $\Z_{p^a}$, of integers modulo a prime power. This result is vastly generalized in Section \ref{sec5} to cover EAQECCs over any finite Galois ring. The paper ends in Section \ref{sec6} with some concluding remarks.



\section{Preliminaries}\label{sec2}  Let $R$ be a finite commutative ring with unity. Let $\text{Hom}(R,\mathbb{C}^*)$ be the set of all additive characters of $R,$ i.e.,  the set of group homomorphisms from $(R,+)$ to $\mathbb{C}^*.$
 A ring $R$ is called a Frobenius ring if there exists an additive character $\chi$ such that  $\text{Hom}(R,\mathbb{C}^*)=R.\chi.$ Any additive character with this property is called a generating character of $R.$ Finite fields, the rings $\mathbb{Z}_N$ of integers modulo $N,$ Galois rings and finite chain rings are a few examples of finite Frobenius rings.
%The following result is well known for finite commutative Frobenius rings.
%
%\begin{prop}\cite{mahmoudi}\label{iso}  The map $\hat{}  : \text{Hom}(R,\mathbb{Q}/\mathbb{Z}) \mapsto\text{Hom}(R,\mathbb{C}^*)$ defined by $\phi \mapsto \hat{\phi}$ with $\hat{\phi}(x)= \exp(2\pi i \phi(x))$ for each $x \in R$ is an isomorphism. Moreover for each $\phi \in \text{Hom}(R,\mathbb{Q}/\mathbb{Z}),$ we have $\text{ker } \phi=\text{ker } \hat{\phi}.$
%\end{prop}
%By the above proposition, there exists a unique $\chi' \in \text{Hom}(\mathcal{R},\mathbb{Q}/\mathbb{Z})$ such that $\chi(x)=\exp(2\pi i \chi'(x))$ for each $x \in \mathcal{R}.$

Throughout this paper, we take $\mathcal{R}$ to be a finite commutative local Frobenius ring with generating character $\chi.$  Further, as $\mathcal{R}$ is a finite commutative local ring, the characteristic, $\text{char}{\mathcal R}$, of $\mathcal{R}$  is a power of a prime number, and the cardinality, $|\mathcal{R}|$, of $\mathcal{R}$ is also a power of prime number. Let  $|\mathcal{R}|=p^a=q,$ and $\text{char}{\mathcal R}=p^b,$ where $p$ is a prime.   Furthermore, with $\zeta = \exp\bigl(\frac{2 \pi i}{p^b}\bigr)$ a primitive $p^b$-th root of unity, we have that $\chi(r)\in \langle \zeta\rangle=\{1,\zeta,\zeta^2,\ldots,\zeta^{p^b-1}\}$ for each $r \in \mathcal{R}.$ 
 
A subset $\mathcal{C}$ of $\mathcal{R}^{n}$ is called an additive code over $\mathcal{R}$ of length $n$ if $\mathcal{C}$ is an additive subgroup of $\mathcal{R}^{n}.$ Clearly, $\mathcal{C} \subseteq \mathcal{R}^{n}$ is an additive code if and only if $\mathcal{C} \subseteq \mathcal{R}^{n}$ is a module over $\mathbb{Z}_{p^b}.$ Since $\mathbb{Z}_{p^b}$ is also a commutative local ring, all minimal (with respect of inclusion) generating sets of finitely generated $\mathbb{Z}_{p^b}$-modules have the same cardinality. In particular, all minimal generating sets of an additive code $\mathcal{C} \subseteq \mathcal{R}^{n}$ have the same cardinality, and this number is called the rank of $\mathcal{C}$, denoted by $\rank(\mathcal{C}).$

In this paper, we will exclusively be concerned with additive codes of even blocklength. In this context, whenever we consider a tuple $(x,y) \in \mathcal{R}^{2n}$, it is to be implicitly understood that $x$ and $y$ both belong to $\mathcal{R}^n$.

 The symplectic weight of a vector $(x,y)=(x_1,x_2,\ldots,x_n,y_1,y_2,\ldots,y_n) \in \mathcal{R}^{2n}$, denoted by $\text{wt}_{\text{s}}(x,y),$ is defined as $$\text{wt}_{\text{s}}(x,y)=|\{ i : (x_i,y_i) \neq 0\}|.$$ The minimum symplectic distance $d_{\mathrm{s}}(A)$ of a subset $A$ of $\mathcal{R}^{2n}$ is defined as  
 $$d_{\mathrm{s}}(A)=\min\{\text{wt}_{\text{s}}(a) :a \in A\setminus \{0\}\}.$$
The symplectic inner product on $\mathcal{R}^{2n}$ is defined as $\langle (a,b) | (a',b')\rangle_{\text{s}} := ba'-b'a$  for $a,b,a',b' \in \mathcal{R}^n.$ (Here $ba'$ and $b'a$ are the dot products in $\mathcal{R}^n.$)

\begin{defin} For an additive code $\mathcal{C} \subseteq \mathcal{R}^{2n}$, we define
\begin{itemize}
\item the symplectic dual, $\mathcal{C}^{\perp_\text{s}}$, of $\mathcal{C}$ as
$$\mathcal{C}^{\perp_{\mathrm{s}}}=\{v \in \mathcal{R}^{2n} :  \langle c | v\rangle_{\text{s}} = 0 \text{~for all~} c \in \mathcal{C}\},$$
% \item the symplectic $p$-dual, $\mathcal{C}^{\perp_{\mathrm{s}}/p}$, of $\mathcal{C}$ as
% $$\mathcal{C}^{\perp_\text{s}/p}=\{v \in \mathcal{R}^{2n} : \langle c | v\rangle_{\text{s}} \in \langle p \rangle \text{~for all~} c \in \mathcal{C}\};$$
\item and the $\chi$-symplectic dual, $\mathcal{C}^{\perp_\chi}$, of $\mathcal{C}$ as $$\mathcal{C}^{\perp_\chi}=\{v \in \mathcal{R}^{2n} : \chi(\langle c | v\rangle_{\text{s}}) = 1 \text{~for all~} c \in \mathcal{C}\}.$$
%\item and the $\chi$-symplectic $p$-dual, $\mathcal{C}^{\perp_\chi/p}$, of $\mathcal{C}$ as
%$$\mathcal{C}^{\perp_\chi/p} =\{ v \in \mathcal{R}^{2n} : \chi(\langle c | v\rangle_{\text{s}}) \in 
%\langle \zeta^{p} \rangle     \text{~for all~} c \in \mathcal{C}\}$$
\end{itemize}
\end{defin}

Note that for an additive code $\mathcal{C} \subseteq \mathcal{R}^{2n}$, the duals defined above are also additive codes in $\mathcal{R}^{2n}.$

\begin{lem}\label{card}  Let $\mathcal{C} \subseteq \mathcal{R}^{2n}$ be an additive code. Then
 $|\mathcal{C}||\mathcal{C}^{\perp_\chi}|=|\mathcal{R}^{2n}|.$
 \end{lem} 
  \begin{proof} The result follows immediately from Lemma 6 of Nadella and Klappenecker  \cite{nadella}.  \end{proof}

\begin{defin}
\label{def:chi-self-ortho}
An additive code $\mathcal{C} \subseteq \mathcal{R}^{2n}$ is called $\chi$-self-orthogonal if $\mathcal{C} \subseteq \mathcal{C}^{\perp_\chi}$.
\end{defin}

\begin{defin} \label{def:chi-ext}
 Let $\mathcal{C} \subseteq \mathcal{R}^{2n}$ be an additive code.
\begin{itemize}
\item A code $\mathcal{C}' \subseteq \mathcal{R}^{2(n+c)}$ is called a $\chi$-self-orthogonal extension of $\mathcal{C}$ if $\mathcal{C}' \subseteq \mathcal{C}'^{\perp_\chi}$, and $\mathcal{C}$ can be obtained from $\mathcal{C}{'} $ by puncturing $\mathcal{C}{'} $ at the coordinates $n+1,n+2,\ldots,n+c, 2n+1,2n+2,\ldots,2n+2c$. The number $c$ is called the entanglement degree. 
\item A $\chi$-self-orthogonal extension of the code $\mathcal{C}$ with the least entanglement degree among all such extensions is called a minimal $\chi$-self-orthogonal extension of $\mathcal{C}.$ The entanglement degree of a minimal $\chi$-self-orthogonal extension of $\mathcal{C}$ is called the minimum entanglement degree of $\mathcal{C}.$
%\item The code $\mathcal{C}'$ is called as a  minimal $\chi$-self-orthogonal extension of $\mathcal{C}$ if $\mathcal{C}'$ is a $\chi$-self-orthogonal extension of the code $\mathcal{C}$ with the minimum entanglement degree.
\end{itemize}
\end{defin}

The reason for the nomenclature of ``entanglement degree'' will get clear when we describe EAQECCs in Sections~\ref{sec:eaqecc} and \ref{sec3}. Briefly, this is the number of entanglement qudits needed in the construction of an EAQECC from an additive code $\mathcal{C}$.


  
  

   
\subsection{Quantum stabilizer codes over local Frobenius rings}\label{sec:stabcodes}

  
Let $$\mathcal{B}=\{| x\rangle : x \in \mathcal{R}\}$$  be an orthonormal basis of ${\mathbb{C}}^{q}.$ The state of a unit system, a qudit, is a superposition of these basis states of the system and is given by

$$| \psi\rangle= \sum\limits_{x \in \mathcal{R}}^{} a_x  | x\rangle, \text{~where~} a_x \in \mathbb{C} \text{~and~}  \sum\limits_{x \in \mathcal{R}}^{} |a_x|^2 =1.$$
An orthonormal basis of the quantum state space of $n$ qudits ${\mathbb{C}}^{q^n}=({\mathbb{C}}^{q})^{\otimes n}$ is given by 

$$\mathcal{B}^{\otimes n}=\{| x\rangle =| x_1\rangle \otimes | x_2\rangle \otimes \cdots \otimes | x_n\rangle  ~:~  x=(x_1,x_2,\ldots,x_n) \in \mathcal{R}^n\}.$$

For $a \in \mathcal{R},$ define two linear maps $X(a)$ and $Z(b)$ on ${\mathbb{C}}^{q}$ by their action on the basis $\mathcal{B}$  as 
$$X(a)(| x\rangle) =| x+a\rangle \text{ and } Z(a)(| x\rangle) =\chi(ax)| x\rangle \text{~for all~} x \in \mathcal{R}.$$
Extend these maps to unitary maps on ${\mathbb{C}}^{q^n}$ as follows:
$$X(a)=X(a_1)\otimes X(a_2)\otimes \cdots\otimes X(a_n) \text{~and~} Z(a)=Z(a_1)\otimes Z(a_2)\otimes \cdots\otimes Z(a_n) \text{~for~} a=(a_1,a_2,\ldots,a_n) \in \mathcal{R}^n.$$
Clearly,  $$X(a)(| x\rangle) =| x+a\rangle \text{ and } Z(a)(| x\rangle) =\chi(ax)| x\rangle \text{~for all~} a, x \in \mathcal{R}^n,$$ where $ax=\sum\limits_{i=1}^{n} a_ix_i$ is the dot product in $\mathcal{R}^n.$
A set $\mathcal{E}_n(\mathcal{R}) := \{X(a)Z(b) : a,b \in \mathcal{R}^n\}$ is called an error basis of the $n$-qudit error group. The following lemma tells us about the multiplicative and commutative properties of elements of the error basis.

\begin{lem}\label{comm}\cite{luerssen} Let $P=X(a)Z(b)$ and $P'=X(a')Z(b')$ with $a,b,a',b' \in \mathcal{R}^n$ be two elements of $\mathcal{E}_n(\mathcal{R}).$ Then $PP'=\chi(ba')X(a+a')Z(b+b').$ Furthermore,  $P$ and $P'$ commute with each other if and only if $\chi(ba'-b'a)=1$ (Here $ba'$ and $b'a$ are the dot products in $\mathcal{R}^n.$)
\end{lem}


To define the $n$-qudit error group, called as the Pauli group, let us fix some notations first.
 Let \begin{equation*} N~=~\left\{\begin{array}{ll}
p^b & \text{ if~} p \text{~is odd} ;\\
2p^b & \text{ if~} p \text{~is even}.\\
\end{array}\right. \end{equation*}
Further, let $\omega \in \mathbb{C}^*$ be a primitive $N$-th root of unity.
\begin{defin} \cite{luerssen}  The Pauli group $\mathcal{P}_n(\mathcal{R})$ is defined as 
$$\mathcal{P}_n(\mathcal{R}) := \{\omega^\ell X(a)Z(b) ~:~ 0 \leq \ell \leq N-1, a, b \in \mathcal{R}^n\}.$$ 
\end{defin}
 Define a map $$\Psi : \mathcal{P}_n(\mathcal{R}) \to \mathcal{R}^{2n} \text{~as~} \Psi(\omega^\ell X(a)Z(b))=(a,b).$$ The map $\Psi$ is a surjective group homomorphism with $\text{ker~}\Psi =\{\omega^\ell I ~:~ 0 \leq\ell \leq N-1 \}.$
 
 The weight of an operator $\omega^\ell X(a)Z(b)=\omega^\ell X(a_1)\otimes X(a_2)\otimes \cdots\otimes X(a_n)Z(a_1)\otimes Z(a_2)\otimes \cdots\otimes Z(a_n) \in \mathcal{P}_n(\mathcal{R})$ is defined as 
 $$\text{wt}(\omega^\ell X(a)Z(b))=|\{ i : (a_i,b_i) \neq 0\}|.$$
 That is, the weight of an operator is defined as the number of non-scalar components of the tensor product that forms the operator.
 
A quantum error correcting code is a $K$-dimensional subspace of $\mathbb{C}^{q^n}.$
A quantum code $\mathcal{Q}$ can detect an error $E \in \mathcal{P}_n(\mathcal{R})$ if and only if  $\langle u | E | v \rangle=\lambda_E\langle u  | v \rangle$ for each $|u \rangle,|v \rangle \in \mathcal{Q}$, where $\lambda_E \in \mathbb{C}$ is a constant depending only on $E$.


 A quantum code $\mathcal{Q}$ has minimum distance $D$ if it can detect all errors of weight at most $D-1$ and cannot detect some error of weight $D$. If $\mathcal{Q} \subseteq \mathbb{C}^{q^n}$ is a quantum code  of dimension $1,$ then clearly, for each  $E \in \mathcal{P}_n(\mathcal{R}),$ we have $\langle u | E | v \rangle=\lambda_E\langle u  | v \rangle$ for each $|u \rangle,|v \rangle \in \mathcal{Q}.$ So in this case, the minimum distance of the quantum code $\mathcal{Q}$ is defined as the largest integer $D$ such that for each non-identity operator $E \in \mathcal{P}_n(\mathcal{R})$ with $\text{wt}(E) < D$, we have $\langle u | E | v \rangle=0$ for all $|u \rangle,|v \rangle \in \mathcal{Q}$. (See, for example, \cite{DG04}.) A quantum code $\mathcal{Q} \subseteq \mathbb{C}^{q^n}$ of dimension $K$ and minimum distance $D$ is referred to as an $((n,K,D))_q$ quantum code. The subscript $q$ may be dropped if there is no ambiguity.

\begin{defin}\cite{luerssen} \begin{itemize}\item[(a)] A subgroup $\mathcal{S}$ of $\mathcal{P}_n(\mathcal{R})$ is called a stabilizer group if 
$\mathcal{S}$ is abelian and $\mathcal{S}  \cap \text{ker } \Psi= \{I_{q{^{n}}}\}.$
\item[(b)] A subspace $\mathcal{Q}$ of $\mathbb{C}^{q^n}$ is called a quantum stabilizer code if there exists a stabilizer group $\mathcal{S}$ such that
$$\mathcal{Q}=\mathcal{Q}(\mathcal{S})~:= \{| x\rangle \in \mathbb{C}^{q^n} ~:~ V | x\rangle =| x\rangle \text{~for all~} V \in \mathcal{S}\}.$$
\end{itemize}
\end{defin}


The following theorem tells us about the dimension of a quantum stabilizer code.
\begin{thm}\label{dim} \cite{luerssen} Let $\mathcal{S} \subseteq \mathcal{P}_n(\mathcal{R}) $ be a stabilizer group and $\mathcal{Q}(\mathcal{S}) \subseteq \mathbb{C}^{q^n}$ be the corresponding quantum stabilizer code. Then, the dimension of $\mathcal{Q}(\mathcal{S})$ is equal to $ q^n/|\mathcal{S}|.$
\end{thm}

%\begin{defin}\cite{ketkar} Let $\mathcal{S} \subseteq \mathcal{P}_n(\mathcal{R}) $ be a stabilizer group and $\mathcal{Q}(\mathcal{S}) \subseteq \mathbb{C}^{q^n}$ be the quantum stabilizer code. Then we have the following:
%\begin{itemize}
%\item[(a)] The quantum code $\mathcal{Q}(\mathcal{S})$ is said to be pure to $t$ if and only if its stabilizer group $\mathcal{S}$ does not contain non-scalar operators of weight less than $t.$
%\item[(b)] The quantum code $\mathcal{Q}(\mathcal{S})$ is said to be pure if and only if it is pure to its minimum distance. By convention, an $((n,1,D))_q$ quantum stabilizer code is assumed to be pure.
%\end{itemize}
%\end{defin}

It is well known that the set of undetectable errors for a quantum stabilizer code $\mathcal{Q}(\mathcal{S})$ are those which commute with all the elements of $\mathcal{S}$ but are not themselves elements of $\mathcal{S}$ (ignoring the overall phase factor of the error and of the elements of $\mathcal{S}$). Hence the error-correction properties of a quantum stabilizer code depend on the centralizer of the stabilizer group. 

\subsection{EAQECCs over local Frobenius rings}\label{sec:eaqecc}

Quantum stabilizer codes can be constructed only from abelian subgroups of $\mathcal{P}_n(\mathcal{R}).$  To construct QECCs from non-abelian subgroups of $\mathcal{P}_n(\mathcal{R})$, there is a framework of entanglement-assisted quantum error correcting codes (EAQECCs), which involves the use of maximally entangled states shared between the transmitter and the receiver. Brun et al.\ \cite{brun} first proposed this construction from non-abelian subgroups of $\mathcal{P}_n(\mathbb{Z}_2)$. They constructed quantum stabilizer codes from stabilizer groups obtained by extending the operators in the non-abelian subgroups into an appropriate higher dimensional space to form abelian subgroups. The qubits added for the extension, called the entanglement qubits, are assumed to be residing error-free at the receiver end.

The basic idea of Brun et al.\ can be extended to construct EAQECCs from non-abelian subgroups of $\mathcal{P}_n(\mathcal{R}).$ For this, we need a method to extend a non-abelian subgroup of $\mathcal{P}_n(\mathcal{R})$ into an appropriate higher dimensional space to form an abelian group.  To this end, we first note that if $\mathcal{S}$ is a subgroup of $\mathcal{P}_n(\mathcal{R}),$ then $\Psi(\mathcal{S}) \subseteq \mathcal{R}^{2n}$ is an additive subgroup of $\mathcal{R}^{2n}.$ Moreover, by Lemma \ref{comm}, we see that an operator $P=\omega^\ell X(a)Z(b) \in \mathcal{P}_n(\mathcal{R})$ commutes with elements of the subgroup $\mathcal{S}$ if and only if $\chi(ba'-b'a)=1$ for each $\omega^{\ell'}X(a')Z(b')\in \mathcal{S}.$ From this, we observe that if $V_1 \in \mathcal{P}_n(\mathcal{R})$ commutes with all elements of $\mathcal{S},$ then $\Psi(V_1) \in \Psi (\mathcal{S})^{\perp_\chi}.$ 

Thus, extending $\mathcal{S}$ to make it an abelian subgroup $\mathcal{S}'$ of $\mathcal{P}_{n+c}(\mathcal{R})$ for some $c$ is equivalent to extending $\mathcal{C} := \Psi(\mathcal{S})$ to $\mathcal{C}' := \Psi(\mathcal{S})' \subseteq  \mathcal{R}^{2(n+c)}$ such that $\mathcal{C}' \subseteq \mathcal{C}'^{\perp_\chi},$ In other words, it is equivalent to constructing a $\chi$-self-orthogonal extension $\mathcal{C}' \subseteq  \mathcal{R}^{2(n+c)}$ of the additive code $\mathcal{C} \subseteq \mathcal{R}^{2n}.$ In Theorem~\ref{tEAC}, we provide a method to construct such a $\chi$-self-orthogonal extension. Then, in Theorem~\ref{thm:eaqecc}, we give a construction of an $((n+c,q^{n+c}/|\mathcal{C}|))$ quantum stabilizer code from $\mathcal{C}'$, which will be the desired EAQECC. As in the Brun et al.\ framework, the $c$ extra qudits involved in the construction are entanglement qudits assumed to be residing error-free at the receiver end. This quantum code is referred to as an $((n,q^{n+c}/|\mathcal{C}|;c))$ EAQECC over $\mathcal{R}$.

As the receiver-end qudits are assumed to be maintained error-free, we note that if $E=\omega^{\ell} X(a,a')Z(b,b') \in \mathcal{P}_{n+c}(\mathcal{R})$, with $a,b \in \mathcal{R}^{n},$ $a',b' \in \mathcal{R}^{c}$, is an error operating on the transmitted codeword, then we must have $a'=b'=(0,0,\ldots,0)\in \mathcal{R}^{c}.$ Thus, only errors of the form $\omega^{\ell} X(a,0)Z(b,0) \in \mathcal{P}_{n+c}(\mathcal{R})$, with $a,b \in \mathcal{R}^{n}$, are assumed to occur in this model. We then say that an $((n,q^{n+c}/|\mathcal{C}| ;c))$ EAQECC has minimum distance $D$ if it can detect all errors of the form $\omega^{\ell} X(a,0)Z(b,0) \in \mathcal{P}_{n+c}(\mathcal{R})$, with $a,b \in \mathcal{R}^{n}$,  of weight at most $D-1$, but it cannot detect some error of this form of weight $D.$  If for an EAQECC $\mathcal{Q}$, all errors of the form $\omega^{\ell} X(a,0)Z(b,0) \in \mathcal{P}_{n+c}(\mathcal{R})$, with $a,b \in \mathcal{R}^{n}$, are detectable, then the minimum distance of the code is defined as the largest integer $D$ such that for each non-identity operator $E= X(a,0)Z(b,0) \in \mathcal{P}_{n+c}(\mathcal{R})$, with $\text{wt}(E) < D$, we have $\langle u | E | v \rangle=0$ for all $|u \rangle,|v \rangle \in \mathcal{Q}$. An $((n,q^{n+c}/|\mathcal{C}| ;c))$ EAQECC with minimum distance $D$ is referred to as an $((n,q^{n+c}/|\mathcal{C}|,D ;c))$ EAQECC over $\mathcal{R}$. 


\subsection{A summary of our main results}
In this paper, we provide a framework to construct EAQECCs from classical additive codes over a finite  commutative local Frobenius ring $\mathcal{R}$. Recall that $|\mathcal{R}| = q = p^a$ and $\text{char}\mathcal{R} = p^b$. Our main result  is as follows:

  \begin{thm}\label{dt1} Let $\mathcal{C} \subseteq \mathcal{R}^{2n}$ be an additive code,  i.e.,  $\mathcal{C}$ is a module over $\mathbb{Z}_{p^b}.$  From $\mathcal{C}$, we can construct an $((n,K, D; c))$ EAQECC over $\mathcal{R},$ where  the number of entanglement qudits needed is $c = \frac12 \rank(\mathcal{C}/(\mathcal{C} \cap \mathcal{C}^{\perp_{{\chi}}}))$, the minimum distance is
\begin{equation*}
D~=~\left\{\begin{array}{ll}
d_{\mathrm{s}}(\mathcal{C}^{\perp_\chi}) & \text{~~if~~} \mathcal{C}^{\perp_\chi} \subseteq \mathcal{C} \\
d_{\mathrm{s}}(\mathcal{C}^{\perp_\chi}\setminus \mathcal{C}) & \text{~~otherwise} \, ,
\end{array}\right. 
\end{equation*}
and the dimension $K$ is bounded as $q^{n+c}/ (|\mathcal{C}| \, {p^{\sum_{t=1}^{b-1} (b-t)\rho_t}}) \leq  K \leq q^{n+c}/|\mathcal{C}|$, the $\rho_t$'s being numbers determined by a certain chain of subcodes of $\mathcal{C}$. Additionally, if $\mathcal{C}$ is free (i.e., it is a free module over $\mathbb{Z}_{p^b}$), then $K = q^{n+c}/|\mathcal{C}|$.
%= \rank(\mathcal{C}/(\mathcal{C} \cap \mathcal{C}^{\perp_\chi}))-\rank(\mathcal{C}/(\mathcal{C} \cap \mathcal{C}^{\perp_\chi/p}))$,  and 
 \end{thm}

Theorem~\ref{dt1} is a direct consequence of Theorems~\ref{tEAC} and \ref{thm:eaqecc} proved in Section~\ref{sec3}. The precise expression for the numbers $\rho_t$ can be found in the restatement of Theorem~\ref{dt1} at the end of that section.

%\begin{cor}\label{cr1}  Let $\mathcal{C} \subseteq \mathcal{R}^{2n}$ be a free additive code,  i.e.,  $\mathcal{C}$ is a free module over $\mathbb{Z}_{p^b}.$  From $\mathcal{C}$, we can construct an $((n,q^{n+c}/|\mathcal{C}|, D; c))$ EAQECC over $\mathcal{R},$ where  the number of entanglement qudits needed is $c = \frac12 \rank(\mathcal{C}/(\mathcal{C} \cap \mathcal{C}^{\perp_{{\chi}}}))  $,  and the minimum distance is
%\begin{equation*}
%D~=~\left\{\begin{array}{ll}
%d_{\mathrm{s}}(\mathcal{C}^{\perp_\chi}) & \text{~~if~~} \mathcal{C}^{\perp_\chi} \subseteq \mathcal{C} \\
%d_{\mathrm{s}}(\mathcal{C}^{\perp_\chi}\setminus \mathcal{C}) & \text{~~otherwise} \,.
%\end{array}\right. 
%\end{equation*}
% \end{cor} 
%

\medskip
 
In the second part of the paper (Sections \ref{sec4}--\ref{sec5}), we explicitly derive the minimum number of pre-shared entanglement qudits required to construct an EAQECC from a submodule over the integer ring $\Z_{p^a}$, and more generally, the minimum number of pre-shared entanglement qudits required to construct an EAQECC from an additive code over a Galois ring.  We state our results in the two theorems below. 

\begin{thm}\label{dt2} Let  $\mathcal{C} \subseteq \mathbb{Z}_{p^a}^{2n}$ be a submodule. From $\mathcal{C}$, we can construct an $((n,K, D; c))$ EAQECC over $\Z_{p^a}$, where the minimum number, $c$, of entanglement qudits needed for the construction is equal to $\frac12 \rank(\mathcal{C}/(\mathcal{C} \cap \mathcal{C}^{\perp_{\mathrm{s}}}))$,  the minimum distance is 
\begin{equation*}
D~=~\left\{\begin{array}{ll}
d_{\mathrm{s}}(\mathcal{C}^{\perp_{\mathrm{s}}}) & \text{~~if~~} \mathcal{C}^{\perp_{\mathrm{s}}} \subseteq \mathcal{C} \\  
d_{\mathrm{s}}(\mathcal{C}^{\perp_{\mathrm{s}}}\setminus \mathcal{C} )  & \text{~~otherwise} \, , \end{array}\right. 
\end{equation*}
and the dimension $K$ is bounded as $p^{a(n+c)}/ (|\mathcal{C}| \, {p^{\sum_{t=1}^{b-1} (b-t)\rho_t}}) \leq  K \leq p^{a(n+c)}/|\mathcal{C}|$, the $\rho_t$'s being numbers determined by a certain chain of subcodes of $\mathcal{C}$. Additionally, if $\mathcal{C}$ is free, then $K = p^{a(n+c)}/|\mathcal{C}|$.
 \end{thm}
%  the dimension $K$ is bounded as $q^{n+c}/ (|\mathcal{C}| {p^{\rho(b-1)}}) \leq  K \leq q^{n+c}/|\mathcal{C}|$ with $\rho=\rank(\mathcal{C}/(\mathcal{C} \cap \mathcal{C}^{\perp_\chi}))-\rank(\mathcal{C}/(\mathcal{C} \cap \mathcal{C}^{\perp_\chi/p}))$,  and 

%\begin{cor}\label{cr2} Let  $\mathcal{C} \subseteq \mathbb{Z}_{p^a}^{2n}$ be a free submodule. From $\mathcal{C}$, we can construct an $((n,p^{a(n+c)}/|\mathcal{C}|, D; c))$ EAQECC over $\Z_{p^a}$, where the minimum number, $c$, of entanglement qudits needed for the construction is equal to $\frac12 \rank(\mathcal{C}/(\mathcal{C} \cap \mathcal{C}^{\perp_{\mathrm{s}}}))  $, and  
%\begin{equation*}D~=~\left\{\begin{array}{ll}
% d_{\mathrm{s}}(\mathcal{C}^{\perp_{\mathrm{s}}}) & \text{~~if~~} \mathcal{C}^{\perp_{\mathrm{s}}} \subseteq \mathcal{C} \,;\\  d_{\mathrm{s}}(\mathcal{C}^{\perp_{\mathrm{s}}}\setminus \mathcal{C} )  & \text{~~otherwise} \,. \end{array}\right. \vspace{-1mm}\end{equation*}
% \end{cor}

\medskip

Theorem~\ref{dt2} follows by putting together the results of Theorems~\ref{tEAC} and \ref{thm:eaqecc}, and Theorem~\ref{4thm}, proved in Sections~\ref{sec3} and \ref{sec4}, respectively. The precise expression for the numbers $\rho_t$ can be found in the restatement of the theorem at the end of Section~\ref{sec4}.


\begin{thm}\label{dt3} Let $\mathcal{C} \subseteq \text{GR}(p^b,m)^{2n}$ be an additive code over the Galois ring $\text{GR}(p^b,m)$.   From $\mathcal{C}$, we can construct an  $((n,K,D;c))$ EAQECC over $\text{GR}(p^b,m),$ where the minimum number, $c$, of entanglement qudits needed for the construction is equal to  $\left\lceil \frac1{2m}\rank(\mathcal{C}/(\mathcal{C} \cap \mathcal{C}^{\perp_{{\Tr}}}))  \right\rceil$, the minimum distance is 
\begin{equation*}
D~=~\left\{\begin{array}{ll}
d_{\mathrm{s}}(\mathcal{C}^{\perp_{{\Tr}}}) & \text{~~if~~} \mathcal{C}^{\perp_{{\Tr}}} \subseteq \mathcal{C} \\
d_{\mathrm{s}}(\mathcal{C}^{\perp_{{\Tr}}}\setminus \mathcal{C} )  & \text{~~otherwise}\, , 
\end{array}\right.
\end{equation*}
and the dimension $K$ is bounded as $p^{bm(n+c)}/ (|\mathcal{C}| \, {p^{\sum_{t=1}^{b-1} (b-t)\rho_t}}) \leq  K \leq p^{bm(n+c)}/|\mathcal{C}|$, the $\rho_t$'s being numbers determined by a certain chain of subcodes of $\mathcal{C}$. Additionally, if $\mathcal{C}$ is free, then $K = p^{bm(n+c)}/|\mathcal{C}|$.
\end{thm}

In the statement of the theorem above, $\mathcal{C}^{\perp_{{\Tr}}}=\{v \in \text{GR}(p^b,m)^{2n} :{\Tr}( \langle v | c\rangle_{\text{s}}) = 0 \; \forall \, c \in \mathcal{C}\}$ is the trace-symplectic dual of $\mathcal{C}$, defined with respect to the generalized trace map ${\Tr}: \text{GR}(p^b,m) \to \mathbb{Z}_{p^b}$. Section~\ref{sec5} contains an exposition of the machinery of the map $\Tr$, and the theorem is a consequence of the results proved in that section. Again, the numbers $\rho_t$ are expressed precisely in the restatement of Theorem~\ref{dt3} at the end of that section.

%\begin{cor}\label{cr3} Let $\mathcal{C} \subseteq \text{GR}(p^b,m)^{2n}$ be a free additive code over the Galois ring $\text{GR}(p^b,m)$.   From $\mathcal{C}$, we can construct an  $((n,p^{bm(n+c)}/{|\mathcal{C}|} ,D;c))$ EAQECC over $\text{GR}(p^b,m),$ where the minimum number, $c$, of entanglement qudits needed for the construction is equal to  $\left\lceil \frac1{2m}\rank(\mathcal{C}/(\mathcal{C} \cap \mathcal{C}^{\perp_{{\Tr}}}))  \right\rceil$, and
%\begin{equation*}D~=~\left\{\begin{array}{ll}
%d_{\mathrm{s}}(\mathcal{C}^{\perp_{{\Tr}}}) & \text{~~if~~} \mathcal{C}^{\perp_{{\Tr}}} \subseteq \mathcal{C} \,;\\ d_{\mathrm{s}}(\mathcal{C}^{\perp_{{\Tr}}}\setminus \mathcal{C} )  & \text{~~otherwise}\,.
%\end{array}\right. \vspace{-1mm}\end{equation*}
%Here, $\mathcal{C}^{\perp_{{\Tr}}}=\{v \in \text{GR}(p^b,m)^{2n} :{\Tr}( \langle v | c\rangle_{\text{s}}) = 0 \; \forall \, c \in \mathcal{C}\}$ is the trace-symplectic dual of $\mathcal{C}$ defined with respect to the generalized trace map ${\Tr}: \text{GR}(p^b,m) \to \mathbb{Z}_{p^b}$.
%\end{cor}


Theorem~\ref{dt3} significantly extends the results of Galindo et al.\ \cite{galindo} and Nadkarni and Garani \cite{nadkarni} obtained for EAQECCs derived from codes over finite fields. For ready reference, we state below the result for finite fields obtained as a corollary of the theorem. This result can also be inferred from the work of Nadkarni and Garani \cite{nadkarni}.

\begin{cor}
\label{cor:Fpm}
 Let $\mathcal{C} \subseteq \mathbb{F}_{p^m}^{2n}$ be an additive code over the finite field $\F_{p^m}$. From $\mathcal{C}$, we can construct an  $((n,p^{m(n+c)}/{|\mathcal{C}|} ,D;c))$ EAQECC over $\mathbb{F}_{p^m},$ where the minimum number, $c$, of entanglement qudits needed for the construction is equal to  
$ \left\lceil  \frac{1}{2m} (\dim_{\F_p}(\mathcal{C}) - \dim_{\F_p}(\mathcal{C} \cap \mathcal{C}^{\perp_{{\tr}}})) \right\rceil$, and
\begin{equation*}
D~=~\left\{\begin{array}{ll}
d_{\mathrm{s}}(\mathcal{C}^{\perp_{\tr}}) & \text{~~if~~}  \mathcal{C}^{\perp_{\tr}} \subseteq \mathcal{C}\\ 
d_{\mathrm{s}}(\mathcal{C}^{\perp_{\tr}}\setminus \mathcal{C} )  & \text{~~otherwise}\,.
\end{array}\right. 
\end{equation*}
Here, $\mathcal{C}^{\perp_{\tr}} := \{v \in \mathbb{F}_{p^m} ^{2n} : \tr( \langle v | c\rangle_{\text{s}}) = 0 \; \forall \, c \in \mathcal{C}\}$ is the trace-symplectic dual of $\mathcal{C}$ defined with respect to the trace map $\tr: \F_{p^m} \to \F_p$ given by $\tr(z) = z+z^p+z^{p^2}+\cdots+z^{p^{m-1}}$.
\end{cor}

  \section{Constructing EAQECCs from Codes over Local Frobenius Rings}\label{sec3}
  
In this section, we provide the details of our method for constructing EAQECCs from additive codes over a finite commutative local Frobenius ring $\mathcal{R}$ with generating character $\chi$. We again recall that $|\mathcal{R}| = q = p^a$ and $\text{char}\mathcal{R} = p^b$. At the heart of the construction is a mechanism to obtain a $\chi$-self-orthogonal extension $\mathcal{C}'$ of an additive code $\mathcal{C}$. This requires some preliminary development.

\subsection{Isotropic Generators and Hyperbolic Pairs}
We start with some definitions.
\begin{defin} \label{def:sympl} A subset $\{a_{11},a_{12},a_{21},a_{22},\ldots,a_{e1},a_{e2}\}$
 of $\mathcal{R}^{2n}$ is said to be a symplectic subset of $\mathcal{R}^{2n}$ if
  $\chi(\langle a_{i1} \mid a_{j1} \rangle_{\text{s}}) = \chi(\langle a_{i2} \mid a_{j2} \rangle_{\text{s}}) = \chi(\langle a_{i1} \mid a_{k 2} \rangle_{\text{s}}) = 1$ and  $\chi(\langle a_{i1} \mid a_{i2} \rangle_{\text{s}})\neq1$ for all $i ,j, k \in \{1,2,\ldots,e\}$ with $i \neq k.$
 \end{defin}

\begin{defin} Let $\mathcal{C} \subseteq \mathcal{R}^{2n}$ be an additive code, i.e.,  $\mathcal{C} \subseteq \mathcal{R}^{2n}$ is a module over $\mathbb{Z}_{p^b}.$ Further, let $\mathcal{G}$ be a generating set of  $\mathcal{C} \subseteq \mathcal{R}^{2n}$ as a $\mathbb{Z}_{p^b}$-module.
\begin{itemize}
\item A generator $g \in \mathcal{G}$ is called an isotropic generator if  $\chi(\langle g | h \rangle_{\text{s}})=1$ for all $h \in \mathcal{G}.$
\item Two generators $g,g' \in \mathcal{G}$ are called a hyperbolic pair if $\chi(\langle g | g' \rangle_{\text{s}})\neq1$, $\chi(\langle g| h \rangle_{\text{s}})= 1$ for all $h \in \mathcal{G} \setminus \{g'\}$, and $\chi(\langle g'| h)\rangle_{\text{s}})=1$ for all $h \in \mathcal{G} \setminus \{g\}$.
\end{itemize}  
\end{defin}
 
Note that a generator is isotropic if and only if  it belongs to $\mathcal{C}^{\perp_\chi}$. Thus, an additive code $\mathcal{C} \subseteq \mathcal{R}^{2n}$ with a generating set $\mathcal{G}$ is $\chi$-self-orthogonal if and only if all the generators in $\mathcal{G}$ are isotropic.
 
 \begin{prop} \label{prop:ext}
 Let $\mathcal{C} \subseteq \mathcal{R}^{2n}$ be an additive code with a generating set containing $e$ hyperbolic pairs, $u_{i1},u_{i2}$, $i = 1,2,\ldots,e$, with the remaining generators being isotropic. Then, $\mathcal{C}$ has a $\chi$-self-orthogonal extension $\mathcal{C}'  \subseteq \mathcal{R}^{2(n+c)}$ if and only if there is a symplectic subset $\{a_{11},a_{12},a_{21},a_{22},\ldots,a_{e1},a_{e2}\} \subset \mathcal{R}^{2c}$ of cardinality $2e$ such that $\chi(\langle u_{i1} \mid u_{i2} \rangle_{\text{s}}) = \chi(\langle a_{i1} \mid a_{i2} \rangle_{\text{s}})$ for $i=1,2,\ldots,e$.
 \end{prop}

 \begin{proof} Let $\mathcal{G} = \{u_{11},u_{12},\ldots,u_{e1},u_{e2},z_1,\ldots,z_d\}$ be a generating set of $\mathcal{C}$, with $z_1,\ldots,z_d$ being the isotropic generators. As the generators live in $\mathcal{R}^{2n}$, we can write them as
 $$
 u_{i1} = (v_i,w_i) \text{ and }  u_{i2} = (x_i,y_i) \text{ for } i = 1,2,\ldots,e,\ \text{ and }  z_j = (v_{e+j},w_{e+j}) \text{ for } j=1,2\ldots,d,
 $$
 where each of the components $u_i,v_i,x_i,y_i$ lies in $\mathcal{R}^n$. 
 
 Now, suppose that $\{a_{11},a_{12},a_{21},a_{22},\ldots,a_{e1},a_{e2}\} \subset \mathcal{R}^{2c}$ is a symplectic subset such that $\chi(\langle u_{i1} \mid u_{i2} \rangle_{\text{s}}) = \chi(\langle a_{i1} \mid a_{i2} \rangle_{\text{s}})$ for $i=1,2,\ldots,e$. Let $a_{i1} = (b_i,c_i)$ and $a_{i2} = (r_i,s_i)$, the components being from $\mathcal{R}^c$. We then extend the components $v_i,w_i,x_i,y_i$ of the generators in $\mathcal{G}$ to $v_i',w_i',x_i',y_i' \in \mathcal{R}^{n+c}$ as follows:
 $$
 v_i' = (v_i,-b_i), \ w_i' = (w_i,c_i), \ x_i' = (x_i,-r_i),\ y_i' = (y_i,s_i) \ \text{ for } i = 1,2,\ldots,e,
 $$
 and $v_i' = (v_i,0)$, $w_i' = (w_i,0)$ for $i = e+1,\ldots,e+d$, where $0$ is the zero element of $\mathcal{R}^c$. Finally, set 
 $$
 u_{i1}' := (v_i',w_i') \text{ and }  u'_{i2} = (x'_i,y'_i) \text{ for } i = 1,2,\ldots,e,\ \text{ and }  z'_j = (v'_{e+j},w'_{e+j}) \text{ for } j=1,2\ldots,d.
 $$
 Then, $\mathcal{G}' = \{u'_{11},u'_{12},\ldots,u'_{e1},u'_{e2},z'_1,\ldots,z'_d\}$ generates an additive code $\mathcal{C}' \subseteq \mathcal{R}^{2(n+c)}$, and the generators in $\mathcal{G}'$ are all isotropic. For instance, $\langle u'_{i1} \mid u'_{i2}\rangle_{\text{s}} = v_i'y_i'-w_i'x_i' = v_iy_i - b_is_i - (w_ix_i-c_ir_i) = \langle u_{i1} \mid u_{i2}\rangle_{\text{s}} - \langle a_{i1} \mid a_{i2}\rangle_{\text{s}}$, so that 
 $$
 \chi(\langle u'_{i1} \mid u'_{i2}\rangle_{\text{s}}) = \chi(\langle u_{i1} \mid u_{i2}\rangle_{\text{s}}) \cdot \bigl(\chi(\langle a_{i1} \mid a_{i2}\rangle_{\text{s}})\bigr)^{-1} = 1.
 $$
It follows that $\mathcal{C}'$ is a $\chi$-self-orthogonal extension of $\mathcal{C}$.

Conversely, if $\mathcal{C}'$ is a $\chi$-self-orthogonal extension of $\mathcal{C}$, then it has codewords $u'_{11},u'_{12},\ldots,u'_{e1},u'_{e2}$ that are extensions of the generators $u_{11},u_{12},\ldots,u_{e1},u_{e2}$ that form the hyperbolic pairs in $\mathcal{G}$. That is, we can write
 $$
 u'_{i1} = \bigl((v_i,\hat{v}_i),(w_i,\hat{w}_i)\bigr) \text{ and }  u'_{i2} = \bigl((x_i,\hat{x}_i),(y_i,\hat{y}_i)\bigr) \text{ for } i = 1,2,\ldots,e,
 $$
for some $\hat{u}_i,\hat{v}_i,\hat{x}_i,\hat{y}_i \in \mathcal{R}^c$. Since $\mathcal{C}'$ is $\chi$-self-orthogonal, we have $\chi(\langle u'_{ik} \mid u'_{j\ell}\rangle_{\text{s}}) = 1$ for all $i,j \in \{1,2,\ldots,e\}$ and $k, \ell \in \{1,2\}$. Hence, setting
$$
a_{i1} = (-\hat{v}_i,\hat{w}_i) \text{ and } a_{i2} = (-\hat{x}_i,\hat{y}_i) \text{ for } i = 1,2,\ldots,e,
$$
it is easy to verify that  $\{a_{11},a_{12},a_{21},a_{22},\ldots,a_{e1},a_{e2}\} \subset \mathcal{R}^{2c}$ is a symplectic subset such that $\chi(\langle u_{i1} \mid u_{i2} \rangle_{\text{s}}) = \chi(\langle a_{i1} \mid a_{i2} \rangle_{\text{s}})$ for $i=1,2,\ldots,e$. 
 \end{proof}

From the above proof, we see that a $\chi$-self-orthogonal extension $\mathcal{C}' \subseteq \mathcal{R}^{2(n+c)}$ of an additive code $\mathcal{C} \subseteq \mathcal{R}^{2n}$ can be constructed in two steps:
\begin{itemize}
\item[(1)] Find a generating set $\mathcal{G}$ of the additive code $\mathcal{C} \subseteq \mathcal{R}^{2n}$ as a $\mathbb{Z}_{p^b}$-module, consisting of a set of hyperbolic pairs of generators and a set of isotropic generators.
\item[(2)] Find a symplectic subset of $\mathcal{R}^{2c}$, for some suitable choice of $c$, satisfying the property required by Proposition~\ref{prop:ext}. The desired code $\mathcal{C}'$ can be generated by the set $\mathcal{G}'$ obtained by extending the generators in $\mathcal{G}$ by $2c$ coordinates as prescribed in the proof of the proposition.
\end{itemize}   

In Theorems~\ref{sym} and \ref{tEAC} below, we will show that it is always possible to construct a $\chi$-self-orthogonal extension of an additive code $\mathcal{C} \subseteq \mathcal{R}^{2n}$ by following this two-step prescription. But before stating and proving these results, we first give a useful structural property of any generating set of $\mathcal{C}$ that consists purely of isotropic generators and hyperbolic pairs, as required in the first step above.

Recall that for each $r \in \mathcal{R}$, we have $\chi(r)\in \langle \zeta \rangle=\{1,\zeta,\zeta^2,\ldots,\zeta^{p^b-1}\}$, where $\zeta = \exp\bigl(\frac{2 \pi i}{p^b}\bigr)$ is a primitive $p^b$-th root of unity. The subgroups of $\langle \zeta \rangle$ are $H_t := \langle \zeta^{p^{b-t}} \rangle = \{1,\zeta^{p^{b-t}},\zeta^{2p^{b-t}},\ldots,\zeta^{(p^{t}-1)p^{b-t}}\}$, for $t = 0,1,\ldots,b$.  The subgroup $H_0$ is simply $\langle 1 \rangle = \{1\}$, and the indexing is chosen so that $H_0 \le H_1 \leq \cdots \le H_b = \langle \zeta \rangle$. For an additive code $\mathcal{C} \subseteq \mathcal{R}^{2n}$ and $t = 0,1,\ldots,b$, define
\begin{equation}
\mathcal{C}^{\perp_{\chi,t}} = \{(v',w') \in \mathcal{R}^{2n}: \chi(\langle(v',w') \mid (v,w) \rangle_{\text{s}}) \in H_t \text{ for all } (v,w) \in \mathcal{C}\}.
\label{def:perp_t}
\end{equation}
Note that $\mathcal{C}^{\perp_{\chi,t}}$ is also an additive code. Clearly, we have the nested sequence $\mathcal{C}^{\perp_\chi} = \mathcal{C}^{\perp_{\chi,0}} \subseteq \mathcal{C}^{\perp_{\chi,1}} \subseteq \cdots \subseteq \mathcal{C}^{\perp_{\chi,b}} = \mathcal{R}^{2n}$.


\begin{prop}\label{prop:count} 
Let $\mathcal{C} \subseteq \mathcal{R}^{2n}$ be an additive code, and $\mathcal{G}$ a generating set of $\mathcal{C}$ as a $\mathbb{Z}_{p^b}$-module consisting only of isotropic generators and hyperbolic pairs. Let $\mathcal{G}_{\mathrm{hyp}}$ denote the subset of generators that form the hyperbolic pairs in $\mathcal{G}$. Then, for $t = 0,1,\ldots,b$, we have $\left|\mathcal{G}_{\mathrm{hyp}} \cap \mathcal{C}^{\perp_{\chi,t}}\right| = \rank(\mathcal{C} / (\mathcal{C} \cap \mathcal{C}^{\perp_{\chi}})) - \rank(\mathcal{C} / (\mathcal{C} \cap \mathcal{C}^{\perp_{\chi,t}}))$.
% let $\mathcal{H}$ be the subset of $\mathcal{G}$ consisting of all generators that do not belong to $\mathcal{C}^{\perp_{\chi,t}}$. Then, $|\mathcal{H}| = \rank(\mathcal{C} / (\mathcal{C} \cap \mathcal{C}^{\perp_{\chi,t}}))$.
% $(v,w)$ and $(x,y)$ such that $\chi(\langle(v,w) \mid (x,y) \rangle_{\text{s}}) \notin H$
\end{prop}

%\textcolor{red}{\begin{prop}\label{prop:count} Let $\mathcal{C} \subseteq \mathcal{R}^{2n}$ be an additive code and $\mathcal{G} = \{(v_1,w_1),(v_2,w_2),\ldots, (v_{c+d},w_{c+d}),$ $(x_1,y_1),(x_2,y_2),$ $\ldots, (x_{c},y_{c})\}$ be a minimal generating set of $\mathcal{C} \subseteq \mathcal{R}^{2n}$ as a $\mathbb{Z}_{p^b}$-module such that $(v_i,w_i)$ and $(x_i,y_i)$, $i=1,2,\ldots,c$, are hyperbolic pairs, and $(v_i,w_i)$, $i=c+1,c+2,\ldots,c+d$, are isotropic generators. Further, let  $\mathcal{M}=\{ i ~| ~1 \leq i \leq c,  (x_i,y_i) \in \mathcal{C}^{\perp_{\chi_1}} \}$ and  $\mathcal{N}=\{ j ~| ~1 \leq j \leq c,  (x_j,y_j) \not\in \mathcal{C}^{\perp_{\chi_1}} \}.$ Then $|\mathcal{M}|=\frac12 (\rank(\mathcal{C}/(\mathcal{C} \cap \mathcal{C}^{\perp_\chi}))-\rank(\mathcal{C}/(\mathcal{C} \cap \mathcal{C}^{\perp_{\chi_1}})))$ and $|\mathcal{N}|=\frac12\rank(\mathcal{C}/(\mathcal{C} \cap \mathcal{C}^{\perp_{\chi_1}})).$
%\end{prop} } 
 \begin{proof}
Let $t$ be any integer in $\{0,1,\ldots,b\}$, and set $H = H_t$. Let $\mathcal{H} = \mathcal{G}_{\text{hyp}} \setminus \mathcal{C}^{\perp_{\chi,t}}$ be the set of hyperbolic generators\footnote{Of course, $\mathcal{H}$ also depends on $t$, but we suppress this dependence from the notation for simplicity.} that do \emph{not} belong to $\mathcal{C}^{\perp_{\chi,t}}$. We will first show that $|\mathcal{H}| = \rank(\mathcal{C} / (\mathcal{C} \cap \mathcal{C}^{\perp_{\chi,t}}))$.
%All isotropic generators in $\mathcal{G}$ are in $\mathcal{C}^{\perp_\chi}$, which is always contained in $\mathcal{C}^{\perp_{\chi,t}}$. Hence, isotropic generators cannot be in $\mathcal{H}$.}

Let $(v_i,w_i)$ and $(x_i,y_i)$, $i=1,2,\ldots,c$, constitute the hyperbolic pairs in $\mathcal{G}$. Observe that if $\chi( \langle (v_i,w_i) \mid (x_i,y_i) \rangle_{\text{s}}) \in H$, then by the way hyperbolic pairs are defined, it follows that both $(v_i,w_i)$ and $(x_i,y_i)$ are in $\mathcal{C}^{\perp_{\chi,t}}$. So, the generators constituting such hyperbolic pairs cannot be in $\mathcal{H}$. On the other hand, if $\chi( \langle (v_i,w_i) \mid (x_i,y_i) \rangle_{\text{s}}) \notin H$, then $(v_i,w_i)$ and $(x_i,y_i)$ are both in $\mathcal{H}$. Thus, we may assume that $\mathcal{H}$ consists of the hyperbolic pairs $(v_i,w_i)$ and $(x_i,y_i)$, $i=1,2,\ldots,h$, and hence, $\chi( \langle (v_i,w_i) \mid (x_i,y_i) \rangle_{\text{s}}) \notin H$ for these pairs.

Let $\pi : \mathcal{C}  \to \mathcal{C}/(\mathcal{C} \cap \mathcal{C}^{\perp_{\chi,t}})$ be the canonical projection map that takes $(v,w) \in \mathcal{C}$ to the coset $[v,w] := (v,w) + (\mathcal{C} \cap \mathcal{C}^{\perp_{\chi,t}})$. Since $\pi$ maps any generator in $\mathcal{G} \setminus \mathcal{H}$ to $\mathcal{C} \cap \mathcal{C}^{\perp_{\chi,t}}$, i.e., to the coset $[0,0]$, it must be the case that $\pi(\mathcal{H}) = \{\pi((v_1,w_1)), \pi((x_1,y_1)), \ldots, \pi((v_h,w_h)), \pi((x_h,y_h))\}$ is a generating set for the quotient $\mathcal{C} / (\mathcal{C} \cap \mathcal{C}^{\perp_{\chi,t}})$, again viewed as a $\mathbb{Z}_{p^b}$-module. We claim that for no generator in $\mathcal{H}$ is it the case that its image under $\pi$ is expressible as a $\mathbb{Z}_{p^b}$-linear combination of the $\pi$-images of other generators in $\mathcal{H}$. It then follows that the restriction of $\pi$ to $\mathcal{H}$ is one-to-one, so that $|\mathcal{H}| = |\pi(\mathcal{H})|$, and that $\pi(\mathcal{H})$ is a minimal generating set of $\mathcal{C} / (\mathcal{C} \cap \mathcal{C}^{\perp_{\chi,t}})$, so that $|\pi(\mathcal{H})| = \rank(\mathcal{C} / (\mathcal{C} \cap \mathcal{C}^{\perp_{\chi,t}}))$. Thus, proving the claim suffices to prove that $|\mathcal{H}| = \rank(\mathcal{C} / (\mathcal{C} \cap \mathcal{C}^{\perp_{\chi,t}}))$.

We prove the claim by way of contradiction. Assume, to the contrary, that $\pi((v_1,w_1))$ is expressible as 
$$
\pi((v_1,w_1)) = \eta_1 \cdot \pi((x_1,y_1)) + \sum_{i=2}^h [\theta_i \cdot \pi((v_i,w_i)) + \eta_i \cdot \pi((x_i,y_i))]
$$
for some $\eta_i$'s and $\theta_i$'s from $\mathbb{Z}_{p^b}$. Then, $\pi\bigl((v_1,w_1) - \eta_1 \cdot (x_1,y_1) - \sum_{i=2}^h [\theta_i \cdot (v_i,w_i) + \eta_i \cdot (x_i,y_i)]\bigr) = [0,0]$, or equivalently, $(v_1,w_1) - \eta_1 \cdot (x_1,y_1) - \sum_{i=2}^h [\theta_i \cdot (v_i,w_i) + \eta_i \cdot (x_i,y_i)] =: (v',w')$ belongs to $\mathcal{C} \cap \mathcal{C}^{\perp_{\chi,t}}$.

Consider the quantity $\kappa := \chi( \langle (v',w') \mid (x_1,y_1) \rangle_{\text{s}})$. Since $(v',w') \in \mathcal{C}^{\perp_{\chi,t}}$ and $(x_1,y_1) \in \mathcal{C}$, we must have $\kappa \in H$. On the other hand, 
since $(v_1,w_1)$ and $(x_1,y_1)$ form a hyperbolic pair,  $\kappa$ reduces to $\chi( \langle (v_1,w_1) \mid (x_1,y_1) \rangle_{\text{s}})$, and moreover, since $(v_1,w_1)$ and $(x_1,y_1)$ are in $\mathcal{H}$, we have $\kappa = \chi( \langle (v_1,w_1) \mid (x_1,y_1) \rangle_{\text{s}}) \notin H$. This is the desired contradiction that proves the claim, and hence, the fact that $|\mathcal{H}| = \rank(\mathcal{C} / (\mathcal{C} \cap \mathcal{C}^{\perp_{\chi,t}}))$.

Now, observe that in the special case of $t = 0$, we have $\mathcal{H}$ being all of $\mathcal{G}_{\mathrm{hyp}}$, and consequently, $\left|\mathcal{G}_{\mathrm{hyp}}\right| = \rank(\mathcal{C} / (\mathcal{C} \cap \mathcal{C}^{\perp_{\chi}}))$. More generally, for an arbitrary subgroup $H$, we have that $\mathcal{G}_{\mathrm{hyp}}$ is the disjoint union of $\mathcal{H}$ and $\mathcal{G}_{\mathrm{hyp}} \cap \mathcal{C}^{\perp_{\chi,t}}$; hence, 
\begin{align*}
\left|\mathcal{G}_{\mathrm{hyp}} \cap \mathcal{C}^{\perp_{\chi,t}}\right| \ &= \ \left|\mathcal{G}_{\mathrm{hyp}}\right| - |\mathcal{H}| \\
&= \ \rank(\mathcal{C} / (\mathcal{C} \cap \mathcal{C}^{\perp_{\chi}})) - \rank(\mathcal{C} / (\mathcal{C} \cap \mathcal{C}^{\perp_{\chi,t}})),
\end{align*}
which completes the proof of the proposition.
\end{proof}

We will make use of this result in the proof of Theorem~\ref{tEAC} that is to follow.

\subsection{Constructing a $\chi$-Self-Orthogonal Extension} \label{sec:chi-self-ortho}
 
The proof of the next theorem describes a method to obtain a generating set of an additive code in $\mathcal{R}^{2n}$ consisting purely of isotropic generators and hyperbolic pairs, fulfilling the requirement in the first step of the two-step procedure for constructing a $\chi$-self-orthogonal extension of an additive code $\mathcal{C}$.  

\begin{thm}\label{sym} Let $\mathcal{C} \subseteq \mathcal{R}^{2n}$ be an additive code. There exists a generating set of $\mathcal{C}$ as a $\mathbb{Z}_{p^b}$-module, that consists only of isotropic generators and hyperbolic pairs. Moreover, the number, $c$, of hyperbolic pairs satisfies $2c=\rank(\mathcal{C}/(\mathcal{C} \cap \mathcal{C}^{\perp_\chi}))$.
\end{thm}  

\begin{proof} 
Let $\pi : \mathcal{C}  \to \mathcal{C}/(\mathcal{C} \cap \mathcal{C}^{\perp_\chi})$ be the canonical projection map that takes $(v,w) \in \mathcal{C}$ to the coset $(v,w) + (\mathcal{C} \cap \mathcal{C}^{\perp_\chi})$. Let $\mathcal{T}=\{(a_1,b_1),(a_2,b_2),\ldots,(a_f,b_f)\} \subseteq \mathcal{C} $ be such that $\pi(\mathcal{T}) := \{\pi((a_1,b_1)),\pi((a_2,b_2)),\ldots,\pi((a_f,b_f))\}$ is a minimal generating set of $\mathcal{C}/(\mathcal{C} \cap \mathcal{C}^{\perp_\chi})$ as a $\mathbb{Z}_{p^b}$-module. Thus, $|\mathcal{T}| = |\pi(\mathcal{T})| = \rank(\mathcal{C}/(\mathcal{C} \cap \mathcal{C}^{\perp_\chi}))$.
Further,  let $\mathcal{S}$ be a minimal generating set of the additive code $\mathcal{C} \cap \mathcal{C}^{\perp_\chi}  \subseteq \mathcal{R}^{2n}$ as a $\mathbb{Z}_{p^b}$-module.
 We assert that  $\mathcal{S} \cup \mathcal{T}$ generates $\mathcal{C}$.
 
To prove this assertion,  for any $(v,w) \in \mathcal{C},$ we see that
$\pi(v,w)=z_1\pi((a_1,b_1))+z_2\pi((a_2,b_2))+\ldots+z_f\pi((a_f,b_f))$ 
for some $z_i \in \mathbb{Z}_{p^b}.$  Since $z_1 \pi((a_1,b_1)) + z_2\pi((a_2,b_2)) + \cdots + z_f \pi((a_f,b_f)) = \pi (z_1 (a_1,b_1) + z_2 (a_2,b_2) + \cdots + z_f (a_f,b_f))$, we have $(v,w)-z_1(a_1,b_1)-z_2(a_2,b_2)-\ldots-z_f((a_f,b_f) \in \mathcal{C} \cap \mathcal{C}^{\perp_\chi},$ which implies that $(v,w) \in \langle \mathcal{S} \cup \mathcal{T} \rangle$. This completes the proof of the assertion.

%We further note that, for any $j \in \{f+1,\ldots, f+g\}$, we have $\chi(b_ia_j-b_ja_i) = 1$ for all $i \in \{1,2,\ldots,f+g\}$, since $(a_j,b_j) \in \mathcal{C}^{\perp_\chi}$.

Since the generators in $\mathcal{S}$ belong to $\mathcal{C}^{\perp_\chi}$, they are all isotropic. We will not tamper with $\mathcal{S}$; instead, we will bring $\mathcal{T}$ into hyperbolic-pair form. For $1 \leq i,j \leq f,$ $\chi(b_ia_j-b_ja_i)$ is a $p^b$-th root of unity, so let $\chi(b_ia_j-b_ja_i)=\zeta^{\ell_{i,j}},$ where $\zeta = \exp(\frac{2\pi i}{p^b})$ and $0\leq \ell_{i,j} < p^b.$ Further, for $1 \leq i,j \leq f,$ we note that $\chi(b_ja_i-b_ia_j)=\chi(b_ia_j-b_ja_i)^{-1},$ which gives $\ell_{i,j} \equiv -\ell_{j,i} \pmod{p^b}$. For each $i \in \{1,2,\ldots,f\}$, we must have $\ell_{i,j} \ne 0$ for some $j \in \{1,\ldots,f\}$; otherwise, $(a_i,b_i)$ is in $\mathcal{C} \cap \mathcal{C}^{\perp_\chi}$, so that $\pi((a_i,b_i))=\mathcal{C} \cap \mathcal{C}^{\perp_\chi},$ contradicting the minimality of $\pi(\mathcal{T})$.

Let $\ell_{t,u}$ with $t,u \in \{1,2, \ldots,f\}$ be such that
$$\gcd(\ell_{t,u},p^b)=\min\{\gcd(\ell_{i,j},p^b) : \ell_{i,j}\neq 0 \text{~and~} 1 \leq i,j \leq f \}.$$ 
Clearly, $\gcd(\ell_{t,u},p^b)$ divides $\gcd(\ell_{i,j},p^b)$ for $1 \leq i,j \leq f.$
 As $\ell_{t,u}\neq 0,$  we have  $\chi(b_ta_u-b_ua_t) \neq 1.$   Swap $(a_1,b_1)$ and $(a_2,b_2)$  with $(a_t,b_t)$ and $(a_u,b_u),$ respectively.
 % We now proceed to replace those generators from the set  $\{(a_3,b_3),(a_4,b_4),\ldots,(a_f,b_f)\}$ which do not commute with $(a_1,b_1)$ and $(a_2,b_2)$.

For $3 \leq i \leq f,$   replace $ (a_i,b_i)$ with 
 $$(a_i{'},b_i{'})=(a_i,b_i) +u_i \, (a_1,b_1) + v_i \, (a_2,b_2),$$ 
 where $u_i$ and $v_i$ are solutions of the linear equations 
$$\ell_{1,2}u_i\equiv\ell_{2,i}(\text{mod~}p^b) \ \ \ \text{ and } \ \ \ \ell_{1,2}v_i\equiv-\ell_{1,i}(\text{mod~}p^b).$$
Such $u_i$ and $v_i$ always exist, since $\gcd(\ell_{1,2},p^b)$ divides $\ell_{2,i}$ and $\ell_{1,i}$. By doing this, we get a new set 
$$\mathcal{T}_1 = \{(a_1,b_1),(a_2,b_2),(a_3{'},b_3{'}),\ldots,(a_f{'},b_f{'})\}$$
such that $\chi(b_2a_1-b_1a_2)\neq1$, $\chi(b_j{'}a_1-b_1a_j{'})=1$, and $\chi(b_j{'}a_2-b_2a_j{'})=1$  for $j \in \{3,4,\ldots,f\}$. In other words, $(a_1,b_1)$ and $(a_2,b_2)$ form a hyperbolic pair. Since $\mathcal{T}$ is recoverable from $\mathcal{T}_1$, we see that $\mathcal{S} \cup \mathcal{T}_1$ also generates $\mathcal{C}$. Moreover, $\pi(\mathcal{T})$ is recoverable from $\pi(\mathcal{T}_1) := \{\pi((a_1,b_1)),\pi((a_2,b_2)),\pi((a_3{'},b_3{'})),\ldots,\pi((a_f{'},b_f{'}))$, so $\pi(\mathcal{T}_1)$ is also a minimal generating set of $\mathcal{C}/(\mathcal{C} \cap \mathcal{C}^{\perp_\chi})$. In particular, $|\mathcal{T}_1| = |\pi(\mathcal{T}_1)| = \rank(\mathcal{C}/(\mathcal{C} \cap \mathcal{C}^{\perp_\chi}))$.

By repeatedly applying the above process to the generators $(a_3{'},b_3{'}),\ldots,(a_f{'},b_f{'}),$ we will eventually obtain a set 
 $$\mathcal{T}' = \{(v_1,w_1),(v_2,w_2),\ldots,(v_{c},w_{c}),(x_1,y_1),(x_2,y_2),\ldots, (x_{c},y_{c})\}$$ 
such that $(v_i,w_i)$ and $(x_i,y_i)$, $i=1,2,\ldots,c$, are hyperbolic pairs, $\mathcal{S} \cup \mathcal{T}'$ generates $\mathcal{C}$, and $\pi(\mathcal{T}')$ is a minimal generating set of $\mathcal{C}/(\mathcal{C} \cap \mathcal{C}^{\perp_\chi})$. In particular, $2c = |\mathcal{T}'| = |\pi(\mathcal{T}')| = \rank(\mathcal{C}/(\mathcal{C} \cap \mathcal{C}^{\perp_\chi}))$.
%
%This also implies that no generator in the set $\{(v_1,w_1),(v_2,w_2),\ldots, $ $  (v_{c},w_{c}),(x_1,y_1),(x_2,y_2),\ldots, (x_{c},y_{c})\}$ can be obtained by taking $\mathbb{Z}_{p^b}$-linear combinations of all the other generators (both isotropic and hyperbolic). 
%Note that it is possible that some of the isotropic generators can be obtained by taking linear combinations of the other generators (again both isotropic and hyperbolic).
%Hence we can find a minimal generating set 
% $\{(v_1,w_1),(v_2,w_2),\ldots, $ $ (v_{c+d},w_{c+d}),(x_1,y_1),(x_2,y_2),\ldots, (x_{c},y_{c})\}$ of $\mathcal{C}$ as a $\mathbb{Z}_{p^b}$-module such that $(v_i,w_i)$ and $(x_i,y_i)$, $i=1,2,\ldots,c$, are hyperbolic pairs, and $(v_{c+i},w_{c+i})$, $i=1,2,\ldots,d$, are isotropic generators.  
%
This completes the proof of the theorem. 
 \end{proof}
 
 \begin{rem} In the construction given in the proof above, while $\mathcal{S}$ and $\mathcal{T}'$ are minimal generating sets of $\mathcal{C} \cap \mathcal{C}^{\perp_\chi}$ and $\mathcal{C}/(\mathcal{C} \cap \mathcal{C}^{\perp_\chi}))$, respectively, it is not guaranteed that $\mathcal{S} \cup \mathcal{T}'$ is also a \emph{minimal} generating set of $\mathcal{C}$. Indeed, it is possible that some of the generators in $\mathcal{S}$ can be obtained by taking linear combinations of other generators in $\mathcal{S} \cup \mathcal{T}'$. Thus, we can only say that $\rank(\mathcal{C}) \le |\mathcal{S} \cup \mathcal{T}'| = \rank(\mathcal{C} \cap \mathcal{C}^{\perp_\chi}) + \rank(\mathcal{C}/(\mathcal{C} \cap \mathcal{C}^{\perp_\chi}))$, which gives us the lower bound $2c \ge \rank(\mathcal{C}) - \rank(\mathcal{C} \cap \mathcal{C}^{\perp_\chi})$.
\end{rem}
  
For the second step in the construction of a $\chi$-self-orthogonal extension of an additive code $\mathcal{C}$, we need a suitable symplectic subset. Such a subset can always be obtained from the  generating set of $\mathcal{C}$ guaranteed by Theorem~\ref{sym}, as explained in the proof of the theorem below.
  
  \begin{thm}\label{tEAC}  Let $\mathcal{C} \subseteq \mathcal{R}^{2n}$ be an additive code,  i.e.,  a module over $\mathbb{Z}_{p^b}.$ Then there exists a $\chi$-self-orthogonal extension $\mathcal{C}{'} \subseteq \mathcal{R}^{2(n+c)}$ of $\mathcal{C}$  with entanglement degree $c$ equal to $\frac12 \rank(\mathcal{C}/(\mathcal{C} \cap \mathcal{C}^{\perp_\chi}))$, whose size can be bounded as  $|\mathcal{C}|  \leq |\mathcal{C}{'}| \leq |\mathcal{C}| \cdot p^{\sum_{t=1}^{b-1} (b-t)  \rho_t}$, with $\rho_t:=\rank(\mathcal{C}/(\mathcal{C} \cap \mathcal{C}^{\perp_{\chi,t-1}}))-\rank(\mathcal{C}/(\mathcal{C} \cap \mathcal{C}^{\perp_{\chi,t}}))$. Moreover, if $\mathcal{C}$ is a free module over $\mathbb{Z}_{p^b},$ then $|\mathcal{C}'| =|\mathcal{C}|$.
  %As a consequence, the minimum entanglement degree of $\mathcal{C}$ is at most $(\rank(\mathcal{C})-\rank(\mathcal{C} \cap\mathcal{C}^{\perp_\chi}))/2.$ 
  \end{thm}

\begin{proof} By Theorem \ref{sym}, there exists a generating set  $$\mathcal{G}=\{(v_1,w_1),(v_2,w_2),\ldots, (v_{c+d},w_{c+d}),(x_1,y_1),(x_2,y_2),\ldots, (x_{c},y_{c})\}$$ of $\mathcal{C}$ as a $\mathbb{Z}_{p^b}$-module such that, for $i = 1,2\ldots,c$, the generators $(v_i,w_i)$ and $(x_i,y_i)$ form hyperbolic pairs, and the generators $(v_i,w_i)$, $i = c+1,\ldots,c+d$, are isotropic. Moreover, $2c = \rank(\mathcal{C}/(\mathcal{C} \cap \mathcal{C}^{\perp_\chi})) .$

For $i = 1,2,\ldots,c$, set $a_{i1} := (-\gamma_i e_i,0)$ and $a_{i2} := (0,e_i)$, where $\gamma_i = \langle(v_i,w_i) \mid (x_i,y_i)\rangle_{\text{s}}$, $e_i = (0,\ldots,0,1,0,$ $\ldots,$ $0) \in \mathcal{R}^c$, the $1$ occurring in the $i$th coordinate, and $0$ is the zero element of $\mathcal{R}^c$. It is then easy to verify that $\{a_{11},a_{12},a_{21},a_{22},\ldots,a_{c1},a_{c2}\} \subset \mathcal{R}^{2c}$ is a symplectic subset such that $\chi(\langle (v_i,w_i) \mid (x_i,y_i) \rangle_{\text{s}}) = \chi(\langle a_{i1} \mid a_{i2} \rangle_{\text{s}})$ for $i=1,2,\ldots,c$.  Hence, by Proposition~\ref{prop:ext}, $\mathcal{C}$ has a $\chi$-self-orthogonal extension $\mathcal{C}' \subseteq \mathcal{R}^{2(n+c)}$. Indeed, $\mathcal{C}'$ is the additive code generated by $\mathcal{G}' = \{u'_{11},u'_{12},\ldots,u'_{c1},u'_{c2},z'_1,\ldots,z'_d\}$, with 
$$
 u_{i1}' := (v_i,-\gamma_i e_i,w_i,0) \text{ and }  u'_{i2} = (x_i,0,y_i,e_i)  \text{ for } i = 1,2,\ldots,c,
 $$
 and  $z'_j = (v_{c+j},0,w_{c+j},0) \text{ for } j=1,2\ldots,d$, where $0$ is the zero element of $\mathcal{R}^c$.

It only remains to prove the assertions about the cardinality of $\mathcal{C}'$. Let $\varphi$ denote the restriction of $\mathcal{C}'$ onto the coordinates of $\mathcal{C}$, i.e., $\varphi$ maps $(v,f,w,g) \in \mathcal{C}'$ to $(v,w) \in \mathcal{C}$. Let $\mathcal{C}_0$ denote the subcode of $\mathcal{C}'$ consisting of codewords of the form $(0,f,0,g)$, for some $f,g \in \mathcal{R}^c$. Then, $\ker \varphi = \mathcal{C}_0$, and hence, $|\mathcal{C}'| = |\mathcal{C}| \cdot |\mathcal{C}_0|$. Since $(0,0,0,0) \in \mathcal{C}_0$, we have $|\mathcal{C}_0| \ge 1$, so that $|\mathcal{C}'| \ge |\mathcal{C}|$. We next work towards an upper bound on $|\mathcal{C}_0|$.

From the definitions of the generating sets $\mathcal{G}$ and $\mathcal{G}'$, we see that $|\mathcal{C}_0|$ equals the number of distinct $(f,g) \in \mathcal{R}^{2c}$ that can be obtained as 
\begin{equation}
\sum_{i=1}^c [\xi_{i,1} \cdot (-\gamma_ie_i,0) + \xi_{i,2} \cdot (0,e_i)] + \sum_{j=1}^d \eta_j \cdot (0,0)
\label{eq:fg}
\end{equation}
for some  $\xi_{i,1},\xi_{i,2},\eta_j \in \mathbb{Z}_{p^b}$ such that 
\begin{equation}
\sum_{i=1}^c [\xi_{i,1} \cdot (v_i,w_i) + \xi_{i,2} \cdot (x_i,y_i)] + \sum_{j=1}^d \eta_j \cdot (v_{c+j},w_{c+j}) \ = \ (0,0).
\label{eq:C0}
\end{equation}
The $\eta_j$ terms do not contribute to the overall sum in \eqref{eq:fg}, so can be dropped from there (but not from \eqref{eq:C0}).

Fix some $\ell \in \{1,2,\ldots,c\}$. Taking the symplectic product with $(x_{\ell},y_{\ell})$ on both sides of \eqref{eq:C0}, and applying the generating character $\chi$, we obtain ${\chi(\gamma_\ell)}^{\xi_{\ell,1}} = 1$, where we recall that $\gamma_\ell = \langle(v_\ell,w_\ell) \mid (x_\ell,y_\ell)\rangle_{\text{s}}$. Similarly, taking the symplectic product with $(v_{\ell},w_{\ell})$ on both sides of \eqref{eq:C0}, and applying the generating character $\chi$, we obtain ${\chi(\gamma_\ell)}^{\xi_{\ell,2}} = 1$. Note that since $(v_{\ell},w_{\ell})$ and $(x_{\ell},y_{\ell})$ form a hyperbolic pair, we have $\chi(\gamma_\ell) \neq 1$. Let $t$ be the least positive integer such that $\chi(\gamma_\ell) \in H_t := \langle \zeta^{p^{b-t}} \rangle$. This means that $\chi(\gamma_\ell)$ is of the form $\exp\bigl(\frac{2\pi i \, z}{p^b}\bigr)$ with $\gcd(z,p^b) = p^{b-t}$. Then, for ${\chi(\gamma_\ell)}^{\xi_{\ell,1}} = {\chi(\gamma_\ell)}^{\xi_{\ell,2}} = 1$, we must have $\xi_{\ell,1}$ and $\xi_{\ell,2}$ being multiples of $p^t$. In particular, if $t=b$, then $\xi_{\ell,1}=\xi_{\ell,2}=0$ (in $\mathbb{Z}_{p^b}$). From all of this, we infer that whenever \eqref{eq:C0} holds, the expression in \eqref{eq:fg} is in fact of the form 
\begin{equation*}
\sum_{t = 1}^{b-1} \sum_{\ell \in \mathcal{I}_t} [\xi_{\ell,1} \cdot (-\gamma_\ell e_\ell,0) + \xi_{\ell,2} \cdot (0,e_\ell)],
\end{equation*}
where $\mathcal{I}_t = \{i : \text{$t$ is the least positive integer such that $\chi(\gamma_i) \in H_t$}\}$, and $\xi_{\ell,1},\xi_{\ell,2} \in \mathbb{Z}_{p^b}$ are multiples of $p^t$. In particular, for $\ell \in \mathcal{I}_t$, we can have at most $p^{b-t}$ possible choices for each of $\xi_{\ell,1}$ and $\xi_{\ell,2}$. Consequently, 
$$
|\mathcal{C}_0| \le \prod_{t=1}^{b-1} \prod_{\ell \in \mathcal{I}_t} p^{2(b-t)} = p^{\sum_{t=1}^{b-1} 2(b-t)|\mathcal{I}_t|}.
$$

Now, observe that $\ell \in \mathcal{I}_t$ iff $(v_{\ell},w_{\ell})$ and $(x_\ell,y_\ell)$ are both in $\mathcal{C}^{\perp_{\chi,t}} \setminus \mathcal{C}^{\perp_{\chi,t-1}}$ --- this can be inferred from the definition of $\mathcal{C}^{\perp_{\chi,t}}$ in \eqref{def:perp_t}. Hence, $2|\mathcal{I}_t| = \left|\mathcal{G}_{\mathrm{hyp}} \cap \mathcal{C}^{\perp_\chi,t} \right| - \left|\mathcal{G}_{\mathrm{hyp}} \cap \mathcal{C}^{\perp_\chi,t-1} \right|$, which from Proposition~\ref{prop:count}, equals $\rho_t := \rank(\mathcal{C}/(\mathcal{C} \cap \mathcal{C}^{\perp_\chi,t-1}))-\rank(\mathcal{C}/(\mathcal{C} \cap \mathcal{C}^{\perp_\chi,t}))$. Thus, we conclude that $|\mathcal{C}_0| \le p^{\sum_{t=1}^{b-1} (b-t) \rho_t}$, and hence, $|\mathcal{C}'| = |\mathcal{C}|\cdot |\mathcal{C}_0|  \le |\mathcal{C}| \cdot p^{\sum_{t=1}^{b-1} (b-t) \rho_t}$.

Finally, if $\mathcal{C}$ is a free module over $\mathbb{Z}_{p^b}$, then \eqref{eq:C0} holds if and only if all the coefficients $\xi_{i,1},\xi_{i,2},\eta_j$ are equal to $0$. Thus, $\mathcal{C}_0$ consists of only the $(0,0,0,0)$ codeword, so that $|\mathcal{C}'|=|\mathcal{C}|\cdot |\mathcal{C}_0| = |\mathcal{C}|$.
\end{proof}

%\begin{proof} By Theorem \ref{sym}, there exists a minimal generating set  $$\{(v_1,w_1),(v_2,w_2),\ldots, (v_{c+d},w_{c+d}),(x_1,y_1),(x_2,y_2),\ldots, (x_{c},y_{c})\}$$ of $\mathcal{C}$ as a $\mathbb{Z}_{p^b}$-module such that, for $i = 1,2\ldots,c$, the generators $(v_i,w_i)$ and $(x_i,y_i)$ form hyperbolic pairs, and the generators $(v_i,w_i)$, $i = c+1,\ldots,c+d$, are isotropic. Moreover, $2c = \rank(\mathcal{C}/(\mathcal{C} \cap \mathcal{C}^{\perp_\chi})) .$
%
%For $i = 1,2,\ldots,c$, set $a_{i1} := (\gamma_i e_i,0)$ and $a_{i2} := (0,e_i)$, where $\gamma_i = \langle(v_i,w_i) \mid (x_i,y_i)\rangle_{\text{s}}$, $e_i = (0,\ldots,0,1,0,\ldots,0) \in \mathcal{R}^c$, the $1$ occurring in the $i$th coordinate, and $0$ is the zero element of $\mathcal{R}^c$. It is then easy to verify that $\{a_{11},a_{12},a_{21},a_{22},\ldots,a_{c1},a_{c2}\} \subset \mathcal{R}^{2c}$ is a symplectic subset such that $\chi(\langle (v_i,w_i) \mid (x_i,y_i) \rangle_{\text{s}}) = \chi(\langle a_{i1} \mid a_{i2} \rangle_{\text{s}})$ for $i=1,2,\ldots,c$. The conclusion of the theorem now follows from Proposition~\ref{prop:ext}.
%\end{proof}

\subsection{Constructing an EAQECC from a $\chi$-Self-Orthogonal Extension}
It only remains to provide a means of constructing an EAQECC from the $\chi$-self-orthogonal extension of $\mathcal{C}$ guaranteed by Theorem~\ref{tEAC}. This is done in the theorem below, which also specifies the dimension and minimum distance of the resulting EAQECC.
 
\begin{thm}\label{thm:eaqecc}
Let $\mathcal{C} \subseteq \mathcal{R}^{2n}$ be an additive code over $\mathcal R$,  and let $\mathcal{C}' \subseteq \mathcal{R}^{2(n+c)}$ be the $\chi$-self-orthogonal extension of $\mathcal{C}$ with entanglement degree $c$, as constructed in Theorem~\ref{tEAC}. Then, there  exists an $((n,q^{n+c}/{|\mathcal{C}'|} ,D;c))$ EAQECC,  where
\begin{equation*}D~=~\left\{\begin{array}{ll}
 d_{\mathrm{s}}(\mathcal{C}^{\perp_\chi}) & \text{~~if~~} \mathcal{C}^{\perp_\chi}\subseteq \mathcal{C}\, ;\\
d_{\mathrm{s}}(\mathcal{C}^{\perp_\chi}\setminus \mathcal{C} )  & \text{~~otherwise} \,.
\end{array}\right. \vspace{-1mm}\end{equation*}
\end{thm}
\begin{proof} From the given $\chi$-self-orthogonal extension, $\mathcal{C'}$, of $\mathcal{C}$, we will first construct a stabilizer group $\mathcal{A} \le \mathcal{P}_{n+c}(\mathcal{R})$ of size equal to $|\mathcal{C}'|$. This is done by adapting the proof of the ``$\Leftarrow$'' part of Theorem~3.12 in \cite{luerssen}.

We start by defining
$$ \mathcal{W}:=\{ \omega^\ell X(v)Z(w) : 0 \leq \ell \leq N-1, (v,w) \in \mathcal{C}'\} \subseteq \mathcal{P}_{n+c}(\mathcal{R}).$$
Note that $\mathcal{W}$ is an abelian subgroup of $\mathcal{P}_{n+c}(\mathcal{R}).$ 
Let $\xi$ be a nontrivial character of $(\mathcal{W},\cdot)$ such that $\xi(\omega^\ell I)=\omega^\ell$ for $0 \leq \ell < N.$ (A character with this property does exist as any character of a subgroup $\langle \omega I \rangle$ of the finite abelian group $\mathcal{W}$  can be
extended to a character of $\mathcal{W}$.) 
%\textcolor{red}{Define a character $\xi$ of $(\mathcal{P}_n(\mathcal{R}),\cdot)$ as $\xi(\omega^\ell X(a)Z(b))=\omega^\ell \text{det}(X(a)Z(b)).$})
Now, define
$$ \mathcal{A}:=\{ \xi(\chi(vw)X(-v)Z(-w))X(v)Z(w) : (v,w) \in \mathcal{C}'\} \subseteq \mathcal{P}_{n+c}(\mathcal{R}).$$
Following the proof of the ``$\Leftarrow$'' part of Theorem~3.12 in \cite{luerssen}, it can be verified that $\mathcal{A}$ is an abelian subgroup of $\mathcal{P}_{n+c}(\mathcal{R})$ such that $\mathcal{A} \cap \text{ker } \Psi= \{I_{q{^{(n+c)}}}\}$. Thus, $\mathcal{A}$ is a stabilizer group of size $|\mathcal{A}| = |\mathcal{C}'|$.

% From this and by using Theorem \ref{tEAC}, we have $|\mathcal{C}| {p^{t(b-1)}} \leq |\mathcal{A}| \leq |\mathcal{C}|$ with $t=\rank(\mathcal{C}/(\mathcal{C} \cap \mathcal{C}^{\perp_\chi}))-\rank(\mathcal{C}/(\mathcal{C} \cap \mathcal{C}^{\perp_{\chi_1}})).$ 


% the last equality being a consequence of the definition of a $\chi$-self-orthogonal extension (Definition~\ref{def:chi-ext}).
 
  By Theorem \ref{dim}, the quantum code $\mathcal{Q}(\mathcal{A}) \subseteq \mathbb{C}^{q^{n+c}}$ has dimension equal to $q^{n+c}/|\mathcal{A}|.$ This proves the existence of an $((n,q^{n+c}/|\mathcal{C}'|; c))$ EAQECC over $\mathcal{R}.$

We now proceed to determine the minimum distance $D$ of the EAQECC. Let $C(\mathcal{A})$  be a centralizer of $\mathcal{A}$ in $\mathcal{P}_{n+c}(\mathcal{R}),$ i.e., $C(\mathcal{A})$ is a set of all the elements of  $\mathcal{P}_{n+c}(\mathcal{R})$ which commute with all the elements of $\mathcal{A}.$ Further, let $E=\omega^{\ell} X(a,a')Z(b,b') \in \mathcal{P}_{n+c}(\mathcal{R})$ with $a,b \in \mathcal{R}^{n},$ $a',b' \in \mathcal{R}^{c}$ be an error occurred on the transmitted codeword. As the receiver-end qudits are assumed to be maintained error-free, we have  $E=\omega^{\ell} X(a,0)Z(b,0) \in \mathcal{P}_{n+c}(\mathcal{R}).$ Further, as the overall phase factor of the error does not matter, we can assume that $E= X(a,0)Z(b,0) \in \mathcal{P}_{n+c}(\mathcal{R})$; indeed, from now on, in this proof, we will ignore the overall phase factor of the elements of $C(\mathcal{A})$ and $\mathcal{A}.$ Note that an element $E= X(a,0)Z(b,0) \in C(\mathcal{A})$  if and only if $(a,b) \in \mathcal{C}^{\perp_\chi}.$

If $E \not\in C(\mathcal{A})$, then we have $EF\neq FE$ for some $F=\omega^{k} X(c,c')Z(d,d') \in \mathcal{A}.$ That is, we have $\chi(b\cdot c-a\cdot d)\neq 1.$ Using this, we observe that $$\langle u | E| v \rangle=\langle u | EF| v \rangle=\chi(b\cdot c-a\cdot d) \langle u | FE| v \rangle=\chi(b\cdot c-a\cdot d) \langle u | E| v \rangle.$$
As $\chi(b\cdot c-a\cdot d)\neq1,$ we get $\langle u | E| v \rangle=0.$ This implies that $E$ is detectable. 
Further, when $E \in\mathcal{A},$ we have $\langle u | E| v \rangle=\langle u |  v \rangle,$ hence $E$ is detectable. Now we assume that $E \in C(\mathcal{A}) \setminus \mathcal{A}$. Here we assert that $E$ is not detectable. We will prove this by contradiction. Suppose that $E$ is detectable. This implies that for each $|v \rangle \in \mathcal{Q}(\mathcal{A}),$ we have $E|v \rangle =\lambda_E |v \rangle$ some complex scalar $\lambda_E.$ Note that for the projection operator $P= \frac{1}{|\mathcal{A}|}\sum\limits_{F \in \mathcal{A}}^{} F,$  we have $EP=\lambda_EP,$ which implies that $\lambda_E \neq 0.$ Now consider an abelian subgroup $\mathcal{A}'$ of $\mathcal{P}_{n+c}(\mathcal{R})$ generated by $\lambda_E^{-1}E$ and all elements of $\mathcal{A}.$ By Theorem \ref{dim}, the quantum code $\mathcal{Q}(\mathcal{A}') \subseteq \mathbb{C}^{q^{n+c}}$ has dimension equal to $q^{n+c}/|\mathcal{A}'|<q^{n+c}/|\mathcal{A}|.$ This implies that not all elements of $\mathcal{Q}(\mathcal{A})$ remain invariant under $\lambda_E^{-1}E.$ This is a contradiction to the detectability of $E.$ This proves our assertion.

Thus an error $E= X(a,0)Z(b,0) \in \mathcal{P}_{n+c}(\mathcal{R})$ is not detectable if and only if $E \in C(\mathcal{A}) \setminus \mathcal{A}$,  i.e., an error $E= X(a,0)Z(b,0) \in \mathcal{P}_{n+c}(\mathcal{R})$ is not detectable if and only if $(a,b) \in \mathcal{C}^{\perp_\chi} \setminus \mathcal{C}.$  From this, we see that if $\mathcal{C}^{\perp_\chi} \setminus \mathcal{C} \not=  \emptyset,$ then $D=d_{\mathrm{s}}(\mathcal{C}^{\perp_\chi}\setminus \mathcal{C}).$ Now, suppose that $\mathcal{C}^{\perp_\chi} \subseteq \mathcal{C}.$  Then, there is no element in  $C(\mathcal{A}) \setminus \mathcal{A}$ of the form  $X(a,0)Z(b,0),$ hence all errors are detectable. Note that  $E=X(a,0)Z(b,0) \in \mathcal{P}_{n+c}(\mathcal{R})$ with $\text{wt}(E) < d_{\mathrm{s}}(\mathcal{C}^{\perp_\chi})$ does not belong to $C(\mathcal{A})$ as $(a,b) \not \in \mathcal{C}^{\perp_\chi} .$ This implies that for each non-identity operator $E=X(a,0)Z(b,0) \in \mathcal{P}_{n+c}(\mathcal{R}),$ with $\text{wt}(E) < d_{\mathrm{s}}(\mathcal{C}^{\perp_\chi}),$ we have $\langle u | E| v \rangle=0$ for each $| u \rangle, | v \rangle \in \mathcal{Q}(\mathcal{A})$.  Further,  for each $(a,b) \in \mathcal{C}^{\perp_\chi}\subseteq \mathcal{C},$ we see that $E =X(a,0)Z(b,0) \in \mathcal{A},$ which implies that $ \langle u | E| v \rangle=\langle u | v \rangle.$ Thus the minimum distance $D$ of the code is equal to $d_{\mathrm{s}}(\mathcal{C}^{\perp_\chi}).$ This completes the proof of the theorem.\end{proof}


%Now we proceed to determine the distance $D$ of the EAQECC. If $\mathcal{C}^{\perp_\chi}=\mathcal{C}$,  then by the assumed upper bound on $c$, we must have $c=0$. Moreoever, by Lemma \ref{card}, we also have $|\mathcal{C}| = q^n$, which implies that the dimension of the quantum code $\mathcal{Q}(\mathcal{A})$ is equal to $q^{n}/|\mathcal{C}|=1$; hence, $\mathcal{Q}(\mathcal{A})$ is a pure code. By the convention for a pure code, the distance $D$ of the code $\mathcal{Q}(\mathcal{A})$  is taken to be the minimum of the weights of the non-scalar operators of  the centralizer of $\mathcal{A}.$ This, in the present case, is equal to $d_{\mathrm{s}}(\mathcal{C}^{\perp_\chi})$; hence, $D=d_{\mathrm{s}}(\mathcal{C}^{\perp_\chi}).$
%
%Now we assume that $\mathcal{C} \neq \mathcal{C}^{\perp_\chi}.$  It is well known that the undetectable errors are those that commute with all the elements of $\mathcal{A}$ but are not the elements of $\mathcal{A}$ (by ignoring the overall phase factor of the error and of the elements of $\mathcal{A}$). Let $E=\omega^{\ell} X(a,a')Z(b,b') \in \mathcal{P}_{n+c}(\mathcal{R})$ with $a,b \in \mathcal{R}^{n},$ $a',b' \in \mathcal{R}^{c}$ be an error occurred on the transmitted codeword. As the receiver-end qudits are assumed to be maintained error-free, we have  $E=\omega^{\ell} X(a,0)Z(b,0) \in \mathcal{P}_{n+c}(\mathcal{R}).$  Further, as $E$ commutes with all elements of $\mathcal{A},$ we get $\chi(b \cdot c+0 \cdot c'-a \cdot d-0 \cdot d')=1$ for each $(c,c',d,d') \in \mathcal{C}' \subseteq \mathcal{R}^{2(n+c)}.$ This gives  $\chi(b\cdot c-a\cdot d)=1$ for each $(c,d) \in \mathcal{C} \subseteq \mathcal{R}^{2n},$ which implies that $(a,b) \in \mathcal{C}^{\perp_\chi}.$ From this and using the fact that undetectable errors are those that commute with all the elements of $\mathcal{A}$ but are not the elements of $\mathcal{A}$ (ignoring overall phase factors), we get $D=d_{\mathrm{s}}(\mathcal{C}^{\perp_\chi}\setminus \mathcal{C} \cap\mathcal{C}^{\perp_\chi}).$ This completes the proof of the theorem.
% \end{proof}

Observe that Theorem~\ref{dt1}, which is restated below with the numbers $\rho_t$ precisely specified, follows readily from Theorems~\ref{tEAC} and \ref{thm:eaqecc}.

\begin{thm}[Restatement of Theorem~\ref{dt1}]  Let $\mathcal{C} \subseteq \mathcal{R}^{2n}$ be an additive code,  i.e.,  $\mathcal{C}$ is a module over $\mathbb{Z}_{p^b}.$  From $\mathcal{C}$, we can construct an $((n,K, D; c))$ EAQECC over $\mathcal{R},$ where the number of entanglement qudits needed is $c = \frac12 \rank(\mathcal{C}/(\mathcal{C} \cap \mathcal{C}^{\perp_{{\chi}}}))$, the minimum distance is
\begin{equation*}
D~=~\left\{\begin{array}{ll}
d_{\mathrm{s}}(\mathcal{C}^{\perp_\chi}) & \text{~~if~~} \mathcal{C}^{\perp_\chi} \subseteq \mathcal{C} \\
d_{\mathrm{s}}(\mathcal{C}^{\perp_\chi}\setminus \mathcal{C}) & \text{~~otherwise} \, ,
\end{array}\right. 
\end{equation*}
and the dimension $K$ is bounded as $q^{n+c}/ (|\mathcal{C}| \, {p^{\sum_{t=1}^{b-1} (b-t)\rho_t}}) \leq  K \leq q^{n+c}/|\mathcal{C}|$, with $\rho_t:=\rank(\mathcal{C}/(\mathcal{C} \cap \mathcal{C}^{\perp_{\chi,t-1}}))-\rank(\mathcal{C}/(\mathcal{C} \cap \mathcal{C}^{\perp_{\chi,t}}))$. Additionally, if $\mathcal{C}$ is free (i.e., it is a free module over $\mathbb{Z}_{p^b}$), then $K = q^{n+c}/|\mathcal{C}|$.
\end{thm}
 
\begin{rem}  We point out here that Lee and Klappenecker \cite[Proposition 2]{lee} provided a method to construct an EAQECC from a free linear code over a finite Frobenius ring. The construction proposed in \cite[Proposition 2]{lee} relies crucially on \cite[Theorem 5]{lee}, which is analogous to our Theorem~\ref{sym}. However, we found a gap in the proof of \cite[Theorem 5]{lee} that could not readily be filled, namely, that replacing $\mathbf{w}_k$ with $\mathbf{w}'_{k-2} = e_{k,i} \mathbf{w}_k - \cdots $ may not result in a basis of $R^{2n}$, as $e_{k,i}$ may not be a unit in the ring $R.$ Our results in this section provide the missing pieces in the construction of EAQECCs over finite commutative \emph{local} Frobenius rings.\end{rem}

Our aim next is to determine the minimum number, $c_{\min}$, of entanglement qudits needed to construct an EAQECC from an additive code $\mathcal{C}$ over the Galois ring, $\text{GR}(p^b,m)$, of characteristic $p^b$ and cardinality $p^{mb}$. To bound $c_{\min}$ from below, we argue as follows: From Theorem~\ref{sym}, $\mathcal{C}$ has a generating set with $c = \frac12 \rank(\mathcal{C}/(\mathcal{C} \cap \mathcal{C}^{\perp_\chi}))$ hyperbolic pairs. If $\mathcal{C}'$ is a minimal $\chi$-self-orthogonal extension of $\mathcal{C}$, with entanglement degree $c_{\min}$, then, by Proposition~\ref{prop:ext}, there is a symplectic subset of $\mathcal{R}^{2c_{\min}}$ of cardinality $2c$. Hence, $2c$ can be bounded above by the size of the largest symplectic subset constructible on $\mathcal{R}^{2c_{\min}}$, which we obtain in terms of $c_{\min}$. This yields a lower bound on $c_{\min}$, as desired. A matching upper bound is obtained by constructing a symplectic subset as stipulated in Proposition~\ref{prop:ext}. This program is carried out for the case of integer rings $\mathbb{Z}_{p^a}$ in the next section, while the treatment for general Galois rings is presented in Section~\ref{sec5}.

\section{EAQECCs from Codes over the Ring $\mathbb{Z}_{p^a}$} \label{sec4}

The map $\chi: \mathbb{Z}_{p^a} \to \mathbb{C}^*$ defined by $\chi(z) = \zeta^z$, with $\zeta=\exp\bigl(\frac{2 \pi i}{p^a}\bigr)$, is a generating character of the ring $\mathbb{Z}_{p^a}$. In particular, $\chi(z) = 1$ iff $z = 0$, and $\chi(z) \in \langle \zeta^{p^t} \rangle $ iff $z \equiv 0 \pmod{p^t}$. Hence, for an additive (linear) code $\mathcal{C} \subseteq \mathbb{Z}_{p^a}^{2n}$, we have $\mathcal{C}^{\perp_\chi}=\mathcal{C}^{\perp_{\mathrm{s}}}$ and $\mathcal{C}^{\perp_{\chi,t}}=\mathcal{C}^{\perp_{\mathrm{s},t}}$, where for $t = 0,1,\ldots,b$, we define
\begin{equation}
\label{def:Cperp_st} 
\mathcal{C}^{\perp_{\mathrm{s},t}} \ = \  \{(v',w') \in \mathbb{Z}_{p^a}^{2n} ~:~ \langle(v',w') \mid (v,w) \rangle_{\mathrm{s}} \equiv 0 \!\!\!\! \pmod{p^{a-t}} \text{ for all } (v,w) \in \mathcal{C}\}.
\end{equation}


Also the definition of a symplectic subset (Definition~\ref{def:sympl}) of $\mathbb{Z}_{p^a}$ reduces to the following:
\begin{defin} \label{def:sympl2} A subset $\{a_{11},a_{12},a_{21},a_{22},\ldots,a_{e1},a_{e2}\}$
 of $\mathbb{Z}_{p^a}^{2n}$ is said to be a symplectic subset if
  $\langle a_{i1} \mid a_{j1} \rangle_{\text{s}} = \langle a_{i2} \mid a_{j2} \rangle_{\text{s}} = \langle a_{i1} \mid a_{k 2} \rangle_{\text{s}} = 0$ and  $\langle a_{i1} \mid a_{i2} \rangle_{\text{s}} \neq 0$ for all $i ,j, k \in \{1,2,\ldots,e\}$ with $i \neq k.$
 \end{defin}

We will first derive an upper bound on the cardinality of a symplectic subset of $\mathbb{Z}_{p^a}^{2n}$, and then, using this bound, we will derive an explicit form of the minimum number of entanglement qudits needed for an EAQECC constructed from an additive (linear) code over $\mathbb{Z}_{p^a}.$   Our upper bound is in fact proved for the more general notion of a quasi-symplectic subset of ${\mathbb{Z}^{2n}_{p^a}}$, defined next. Here, for $z \in \mathbb{Z}_{p^a}$, we use the notation $\overline{z}$ to denote the residue $z \!\! \mod p$. 
\begin{defin} A subset $\{a_{11},a_{12},a_{21},a_{22},\ldots,a_{e1},a_{e2}\}$
 of ${\mathbb{Z}^{2n}_{p^a}}$ is said to be a quasi-symplectic subset if
 \begin{enumerate}
 \item[(a)] $\langle a_{i1}, a_{j1} \rangle_{\text{s}}=\langle a_{i1}, a_{k 2} \rangle_{\text{s}}=0$ for all $i ,j, k  \in \{1,2,\ldots,e\}$ with $i \neq k.$
 \item[(b)] there exists a subset $\mathcal{J}$ of $\{1,2,\ldots,e \}$ such that $\langle a_{i1}, a_{i2} \rangle_{\text{s}}\neq0$ for all $i \not\in \mathcal{J}$, and $\{\overline{a}_{j1} \}_{j \in \mathcal{J}}$ is a linearly independent set over $\mathbb{Z}_p.$
 \end{enumerate}
\end{defin}



%\begin{lem}\label{LemLI} If $\{a_1,a_2,\cdots,a_t\}$ is a subset of $\mathbb{Z}^n_{p^a}$ such that $\{\overline{a}_1,\overline{a}_2,\cdots,\overline{a}_t\}$ is a linearly independent set over $\mathbb{Z}_p,$ then   $\{a_1,a_2,\cdots,a_t\}$ is linearly independent set over $\mathbb{Z}_{p^a}.$ 
%\end{lem}
%\begin{proof} To prove this result, we will apply induction on $a \geq 1.$ Clearly, result holds for $a=1.$  Now we assume that $k\geq 2$ be a fix integer and result holds for integer $a =k-1,$ i.e., we have  if $\{\overline{a}_1,\overline{a}_2,\cdots,\overline{a}_t\}$ is a linearly independent set over $\mathbb{Z}_p,$ then
% $\{a_1,a_2,\cdots,a_t\} \subseteq \mathbb{Z}^n_{p^{k-1}}$ is linearly independent set over $\mathbb{Z}_{p^{k-1}}.$
%
%  Now we will 
%prove that the result holds for $a=k.$ For that, suppose that $ \sum\limits_{\ell=1}^{t} r_\ell a_\ell=0$ in 
%$\mathbb{Z}^n_{p^k}$ with $r_\ell \in \mathbb{Z}_{p^k}$ for each $\ell.$ From this, we get $ 
%\sum\limits_{\ell=1}^{t} \overline{r}_\ell \overline{a}_\ell=0$ in $\mathbb{Z}^n_{p}.$ As 
%$\{ \overline{a}_1,\overline{a}_2,\cdots,\overline{a}_t\}$ is a linearly independent set over 
%$\mathbb{Z}_p,$ we have $\overline{r}_\ell=0$ for all $\ell.$
% This implies that $r_\ell=p {{r'}_\ell}$ for  some 
% ${{r'}_\ell} \in \mathbb{Z}_{p^{k-1}}.$ From this, 
% we get $p( \sum\limits_{\ell=1}^{t} r'_\ell a_\ell)=0$  in $\mathbb{Z}^n_{p^k},$
%  which implies that  $ \sum\limits_{\ell=1}^{t} r'_\ell a_\ell=0$  in $\mathbb{Z}^n_{p^{k-1}}$ (Viewing $a_\ell$ as elements of $\mathbb{Z}^n_{p^{k-1}}.$) Now by induction, we have $\{a_1,a_2,\cdots,a_t\}$ is linearly independent set over $\mathbb{Z}_{p^{k-1}},$ which implies 
%  that $r'_\ell=0$ for each $\ell.$ This implies that $\{a_1,a_2,\cdots,a_t\}$ is linearly independent 
%  set over $\mathbb{Z}_{p^{k}},$  which further implies that the result for $a=k.$ This completes the proof of lemma.
%
%\end{proof}

Note that a symplectic subset satisfies the definition of a quasi-symplectic subset by setting $\mathcal{J} = \emptyset$. In the following theorem, we provide an upper bound on the size of a quasi-symplectic subset of ${\mathbb{Z}^{2n}_{p^a}}.$
\begin{thm} \label{tquasi} If $\{b_{11},b_{12},b_{21},b_{22},\ldots,b_{e1},b_{e2}\}$ is a quasi-symplectic subset of ${\mathbb{Z}^{2n}_{p^a}},$ then $e \leq n.$ 
\end{thm}
\begin{proof} We will apply induction on $a \geq 1$ to prove this result. For $a=1,$ suppose that $\{b_{11},b_{12},b_{21},b_{22},\ldots,b_{e1},b_{e2}\}$ is a quasi-symplectic subset of ${\mathbb{Z}^{2n}_{p}}$. By definition of quasi-symplectic subset, we have $\langle b_{i1}, b_{j1} \rangle_{\text{s}}=0$ and $\langle b_{i1}, b_{k 2} \rangle_{\text{s}}=0$ for all $i ,j, k  \in \{1,2,\ldots,e\}$ with $i \neq k$, and there exists a subset $\mathcal{J}$ of $\{1,2,\ldots,e \}$ such that $\langle b_{t1}, b_{t2} \rangle_{\text{s}}\neq0$ for all $ t \not\in \mathcal{J}$ and $\{b_{v1} \}_{v \in \mathcal{J}}$ is a linearly independent set over $\mathbb{Z}_p$. Here, we assert that $\{b_{11},b_{21},\ldots,b_{e1}\}$ is a linearly independent set over $\mathbb{Z}_p.$ To see this, suppose that $\sum\limits_{\ell=1}^{e}r_\ell b_{\ell 1} =0$ with $r_\ell \in \mathbb{Z}_p.$ From this and using the fact that $\langle b_{k1}, b_{j2}\rangle_{\text{s}} =0$ and $\langle b_{j1}, b_{j2}\rangle_{\text{s}} \neq 0$ for $k \in \{1,2,\ldots,e\}$ and $j \in \mathcal{J}^\mathsf{c}$ with $k \neq j,$ we get $r_t\langle b_{t1}, b_{t2}\rangle_{\text{s}} =0$ for $t \in \mathcal{J}^\mathsf{c}.$ This implies that $r_t=0$ for each $t \in \mathcal{J}^\mathsf{c}.$ So we have $\sum\limits_{ \ell \in \mathcal{J}}^{}  r_\ell b_{\ell 1} =0.$ As $\{b_{\ell 1}\}_{\ell \in \mathcal{J}}$ is a linearly independent set over $\mathbb{Z}_p,$ we get $r_\ell =0$ for each $\ell \in \mathcal{J}.$ This proves the assertion. Now consider a linear code $\mathcal{C}$ over $\mathbb{Z}_p$ generated by  $\{b_{11},b_{21},\ldots,b_{e1}\}.$ Clearly, $\mathcal{C} \subseteq \mathcal{C}^{\perp_{\mathrm{s}}}.$ From this and using Remark 3.8 of  Gluesing-Luerssen and Pllaha  \cite{luerssen}, we get $|{\mathbb{Z}_p}^{2n}|=|\mathcal{C}||\mathcal{C}^{\perp_{\mathrm{s}}}|\geq |{\mathbb{Z}_p}|^{2e}.$ This implies that $e \leq n,$ which completes the proof for $a=1.$ 

Now we assume that $a\geq 2.$ As $\{b_{11},b_{12},b_{21},b_{22},\ldots,b_{e1},b_{e2}\}$
is a quasi-symplectic subset of ${\mathbb{Z}^{2n}_{p^a}},$  we have $\langle b_{i1}, b_{j1} \rangle_{\text{s}}=0$ and $\langle b_{i1}, b_{k 2} \rangle_{\text{s}}=0$ for all $i ,j, k  \in \{1,2,\ldots,e\}$ with $i \neq k$ and there exists a subset $\mathcal{J}$ of $\{1,2,\ldots,e \}$ such that $\langle b_{t1}, b_{t2} \rangle_{\text{s}}\neq0$ for all $ t \not\in \mathcal{J}$ and $\{\overline{b}_{v1} \}_{v \in \mathcal{J}}$ is a linearly independent set over $\mathbb{Z}_p.$ 
 Note that by adding indices from $\mathcal{J}^\mathsf{c}$ to $\mathcal{J}$ if necessary and interchanging $b_{j1}$ and $b_{j2}$ if needed, we can assume that for each $\ell \in \mathcal{J}^\mathsf{c},$
 both $\overline{b}_{j1}$ and $\overline{b}_{j2}$ are linearly dependent on the set $\{\overline{b}_{\ell 1}\}_{\ell \in \mathcal{J}}.$ That is, for each $\ell \in \mathcal{J}^\mathsf{c},$ we have

$$\overline{b}_{\ell 1}=\sum\limits_{ j \in \mathcal{J}}^{}  u_j \overline{b}_{j 1} \text{~and~} \overline{b}_{\ell 2}=\sum\limits_{ j \in \mathcal{J}}^{}  v_j \overline{b}_{j 1} \text{~for ~} u_j, v_j \in \mathbb{Z}_p.$$

This implies that for each $\ell \in \mathcal{J}^\mathsf{c},$ we have

\begin{equation}\label{e1}
b_{\ell 1}=\sum\limits_{ j \in \mathcal{J}}^{}  u_j b_{j 1}+pc_{\ell 1} \text{~and~} {b}_{\ell 2}=\sum\limits_{ j \in \mathcal{J}}^{}  v_j b_{j 1}+pc_{\ell 2} \text{~with~} u_j, v_j \in \mathbb{Z}_p \text{~and~} c_{\ell 1}, c_{\ell 2} \in \mathbb{Z}_{p^{a-1}}^{2n} \end{equation}

Recall that for $t, u \in \mathcal{J}$ and $h,i,j,k,\ell \in \mathcal{J}^\mathsf{c}$ with $i \neq k,$ 
we have $\langle b_{t1}, b_{h1}\rangle_{\text{s}} =0,$  $\langle b_{t1}, b_{h2}\rangle_{\text{s}} =0,$ $\langle b_{h1}, b_{u2}\rangle_{\text{s}}=0,$ $\langle b_{i1}, b_{j1}\rangle_{\text{s}}=0,$ $\langle b_{i1}, b_{k2}\rangle_{\text{s}} =0$ and $ \langle b_{\ell1}, b_{\ell2}\rangle_{\text{s}} \neq 0$ in $\mathbb{Z}_{p^{a}}.$

 From this and using \eqref{e1}, we observe that for $t, u \in \mathcal{J}$ and $h,i,j,k,\ell \in \mathcal{J}^\mathsf{c}$ with $i \neq k,$ we have
 \begin{eqnarray}\label{e2}
 &&    \langle b_{t1}, b_{h1}\rangle_{\text{s}} =p\langle b_{t1},  c_{h1}\rangle_{\text{s}} =0,    \langle b_{t1}, b_{h2}\rangle_{\text{s}} =p\langle b_{t1},  c_{h2}\rangle_{\text{s}} =0, \langle b_{h1}, b_{u2}\rangle_{\text{s}} =p\langle c_{h1},  b_{u2}\rangle_{\text{s}} = 0, 
 \nonumber\\ &&
\langle b_{i1}, b_{j1}\rangle_{\text{s}} =p^2\langle c_{i1},  c_{j1}\rangle_{\text{s}} = 0,   \langle b_{i1}, b_{k2}\rangle_{\text{s}} =p^2\langle c_{i1},  c_{k2}\rangle_{\text{s}} = 0
  \text{~and ~} 
 \langle b_{\ell1}, b_{\ell2}\rangle_{\text{s}} =p^2\langle c_{\ell1},  c_{\ell2}\rangle_{\text{s}} \neq 0  \text{~in~} \mathbb{Z}_{p^{a}}. ~ \end{eqnarray} 
 
  
  
For $a=2,$ by using  \eqref{e2}, we see that $\mathcal{J}^\mathsf{c}=\phi.$ This implies that $\{\overline{a}_{11},\overline{a}_{21},\ldots,\overline{a}_{e1}\}$ is linearly independent over $\mathbb{Z}_{p}.$  Thus  a linear code $\mathcal{C}$ over  $\mathbb{Z}_p$ generated by  $\{b_{11},b_{21},\ldots,b_{e1}\}$ satisfies $\mathcal{C} \subseteq \mathcal{C}^{\perp_{\mathrm{s}}}.$ From this and again by using Remark 3.8 of  Gluesing-Luerssen and Pllaha  \cite{luerssen}, we get $|{\mathbb{Z}_p}^{2n}|=|\mathcal{C}||\mathcal{C}^{\perp_{\mathrm{s}}}|\geq |{\mathbb{Z}_p}|^{2e}.$ This implies that $e \leq n.$  Hence the result holds for $a=2.$   Now we assume that $k\geq 3$  be a fix integer and result holds for integer $a=k-2.$
  We will prove the result for $a=k.$ For that, by using \eqref{e2}, we get $\{ \{\tilde{b}_{j1},\tilde{b}_{j2}\}_{j \in \mathcal{J}},$   $ \{\tilde{c}_{\ell1},\tilde{c}_{\ell2}\}_{\ell \in \mathcal{J}^\mathsf{c}}\}$ is quasi symplectic subset of $\mathbb{Z}_{p^{k-2}}^n,$ where $\tilde{c}\equiv c  (\text{mod~}p^{k-2}).$ So by induction, we must have $e \leq n,$ which implies that the result holds for $a=k.$ This completes the proof of the theorem. 
\end{proof}

Since a symplectic subset of $\mathbb{Z}_{p^a}^{2n}$ is also a quasi-symplectic subset, we have the following corollary.

\begin{cor}\label{cor2} If $\{a_{11},a_{12},a_{21},a_{22},\ldots,a_{e1},a_{e2}\}$ is  a symplectic subset of ${\mathbb{Z}^{2n}_{p^a}},$ then $e \leq n.$ \end{cor}
%\begin{proof} To prove this result, we first note that $\chi(r)=e^{2\pi i r/p^a}$ for each $r \in \mathbb{Z}_{p^a},$ which implies that $\chi(r)=1$ if and only if $r=0.$ This implies that if $\{a_{11},a_{12},a_{21},a_{22},\ldots,a_{e1},a_{e2}\}$
% of ${\mathbb{Z}^{2n}_{p^a}}$ is  a symplectic subset of ${\mathbb{Z}^{2n}_{p^a}},$ then
% $\{a_{11},a_{12},a_{21},a_{22},\ldots,a_{e1},a_{e2}\}$ is also a quasi-symplectic subset of ${\mathbb{Z}^{2n}_{p^a}}.$ By Theorem \ref{tquasi}, we get $e \leq n,$ which completes the proof of the corollary.  
%\end{proof}

From this, we obtain an explicit formula for the minimum entanglement degree of a submodule $\mathcal{C} \subseteq \mathbb{Z}_{p^a}^{2n}.$

\begin{thm}\label{4thm} Any minimal $\chi$-self-orthogonal extension of a submodule $\mathcal{C} \subseteq \mathbb{Z}_{p^a}^{2n}$ has entanglement degree equal to $\frac12 \rank(\mathcal{C}/(\mathcal{C} \cap \mathcal{C}^{\perp_{\mathrm{s}}})).$
\end{thm}
\begin{proof} 
Let $c_{\min}$ be the entanglement degree of a minimal $\chi$-self-orthogonal extension, $\mathcal{C}'$, of $\mathcal{C}$. By Theorem~\ref{dt1}, we have $c_{\min} \le  \frac12 \rank(\mathcal{C}/(\mathcal{C} \cap \mathcal{C}^{\perp_{\mathrm{s}}})) $. On the other hand, from Theorem~\ref{sym}, $\mathcal{C}$ has a generating set with $c = \frac12 \rank(\mathcal{C}/(\mathcal{C} \cap \mathcal{C}^{\perp_{\mathrm{s}}}))$ hyperbolic pairs. Then, by Proposition~\ref{prop:ext}, there is a symplectic subset of $\mathbb{Z}_{p^a}^{2c_{\min}}$ of cardinality $2c$. Hence, by Corollary~\ref{cor2}, we have $c \le c_{\min}$, as desired.
\end{proof}

Observe that Theorem~\ref{dt2}, which we restate below with the numbers $\rho_t$ exactly specified,  follows by combining the results of Theorems~\ref{tEAC}, \ref{thm:eaqecc} and \ref{4thm}, using the fact that $\mathcal{C}^{\perp_\chi}=\mathcal{C}^{\perp_{\mathrm{s}}}$ and  $\mathcal{C}^{\perp_{\chi,t}}=\mathcal{C}^{\perp_{\mathrm{s},t}}$, as defined in \eqref{def:Cperp_st}.

\begin{thm}[Restatement of Theorem~\ref{dt2}] Let  $\mathcal{C} \subseteq \mathbb{Z}_{p^a}^{2n}$ be a submodule. From $\mathcal{C}$, we can construct an $((n,K, D; c))$ EAQECC over $\Z_{p^a}$, where the minimum number, $c$, of entanglement qudits needed for the construction is equal to $\frac12 \rank(\mathcal{C}/(\mathcal{C} \cap \mathcal{C}^{\perp_{\mathrm{s}}}))$,  the minimum distance is 
\begin{equation*}
D~=~\left\{\begin{array}{ll}
d_{\mathrm{s}}(\mathcal{C}^{\perp_{\mathrm{s}}}) & \text{~~if~~} \mathcal{C}^{\perp_{\mathrm{s}}} \subseteq \mathcal{C} \\  
d_{\mathrm{s}}(\mathcal{C}^{\perp_{\mathrm{s}}}\setminus \mathcal{C} )  & \text{~~otherwise} \, , \end{array}\right. 
\end{equation*}
and the dimension $K$ is bounded as $p^{a(n+c)}/ (|\mathcal{C}| \, {p^{\sum_{t=1}^{b-1} (b-t)\rho_t}}) \leq  K \leq p^{a(n+c)}/|\mathcal{C}|$, with $\rho_t:=\rank(\mathcal{C}/(\mathcal{C} \cap \mathcal{C}^{\perp_{\mathrm{s},t-1}}))-\rank(\mathcal{C}/(\mathcal{C} \cap \mathcal{C}^{\perp_{\mathrm{s},t}}))$.  Additionally, if $\mathcal{C}$ is free, then $K = p^{a(n+c)}/|\mathcal{C}|$.
 \end{thm}

\section{EAQECCs from Codes over Galois Rings}\label{sec5}
In this section, we will first derive an upper bound on the cardinality of a symplectic subset of $\text{GR}(p^b,m)^{2n}$ and then using this bound, we will derive an explicit formula for the minimum number of entanglement qudits needed for an EAQECC constructed from an additive code over $\text{GR}(p^b,m)$.  To do this, we start by reviewing some basic properties of Galois rings.

A finite commutative ring with unity is called a Galois ring if all its non-units (or equivalently, all its zero divisors  including $0)$ form an ideal generated by some prime number. Any finite field of characteristic $p$ is a Galois ring as all the non-units form a zero ideal which generated by $p$. Another example of a Galois ring is the ring $\mathbb{Z}_{p^a}$ of integers modulo a prime power $p^a$, in which all the non-units form an ideal generated by $p$.

Let $R$ be a Galois ring, in which the set of all zero-divisors of including $0$ form an ideal generated by a prime number $p$. Then the ideal generated by $p$ is the only maximal ideal of ${R}$ and the characteristic of ${R}$ is $p^b$ for some positive integer $b$. The residue field of ${R}$ is given by ${R} / \langle p \rangle \simeq  \mathbb{F}_{p^m}$ for some positive integer $m$ and  $|{R}|=p^{mb}.$ Moreover, the ring $R$ is isomorphic to the quotient ring $\mathbb{Z}_{p^b}[x]/ \langle h(x) \rangle,$ where  $h(x)\in \mathbb{Z}_{p^b}[x]$ is a monic polynomial of degree $m$ such that its reduction $(h(x) \text{ mod } p) \in \mathbb{F}_{p}[x]$ is irreducible and primitive over $\mathbb{F}_{p}.$ Thus any two Galois rings of  the same characteristic and the same cardinality are isomorphic (see \cite[Chp. 14]{wan}). 

Through out this section, let $\text{GR}(p^b,m)$ denote the Galois ring of characteristic $p^b$ and cardinality $p^{mb}.$  Further, let $h(x)\in \mathbb{Z}_{p^b}[x]$ be a monic polynomial of degree $m$ such that its reduction $(h(x) \text{ mod } p) \in \mathbb{F}_{p}[x]$ is irreducible and primitive over $\mathbb{F}_{p}$ and $h(x)$ divides $x^{p^m-1}-1$ in $\mathbb{Z}_{p^b}[x].$  Then we have  $\text{GR}(p^b,m) \cong \mathbb{Z}_{p^b}[x]/ \langle h(x) \rangle.$ If $\theta=x+\langle h(x) \rangle,$ then clearly $h(\theta)=0$ and every element of $\text{GR}(p^b,m) \cong \mathbb{Z}_{p^b}[x]/ \langle h(x) \rangle$ has a unique representation as $$r=r_0+r_1\theta+\cdots+r_{m-1}\theta^{m-1} \text{~with~} r_0,r_1,\ldots,r_{m-1} \in \mathbb{Z}_{p^b}.$$
This implies that $\text{GR}(p^b,m)$ is a free module of rank $m$ over  $\mathbb{Z}_{p^b}$ with $\{1,\theta,\theta^2,\ldots,\theta^{m-1}\}$ as a basis.
 Moreover,   there exists an element $\beta$  in $\text{GR}(p^b,m)$ having the multiplicative order as  $p^{m}-1,$ which is a root of $h(x).$ The set $\mathcal{T}= \{ 0,1, \beta,\ldots,\beta^{p^{m}-2} \}$ is called the Teichm\"{u}ller set of $\text{GR}(p^b,m)$ and each element $z \in \text{GR}(p^b,m)$ can be uniquely expressed as $z=z_0+ z_1 p+ z_2 p^2+\cdots + z_{b-1} p^{b-1},$ where $z_0,z_1,\ldots,z_{b-1} \in \mathcal{T}$ (see \cite[Chp. 14]{wan}). 

 
The generalized Frobenius map $f$ of the Galois ring $\text{GR}(p^b,m)$ to $\mathbb{Z}_{p^b}$ is defined by $$f(z)=z^p_0+ z^p_1 p+ z^p_2 p^2+\cdots + z^p_{b-1} p^{b-1}$$ for each $z=z_0+ z_1 p+ z_2 p^2+\cdots + z_{b-1} p^{b-1} \in \text{GR}(p^b,m),$ where $z_0,z_1,\ldots,z_{b-1} \in \mathcal{T}.$  

For $b = 1,$ the map $f$ reduces to the usual Frobenius automorphism on $\mathbb{F}_{p^m}$ defined by $f(z)=z^p$ for all $z \in \mathbb{F}_{p^m}.$ 
The generalized trace map ${\Tr}$ from $\text{GR}(p^b,m)$ to $\mathbb{Z}_{p^b}$ is defined by 
$${\Tr}(z)=z+f(z)+f^2(z)+\cdots+f^{m-1}(z) \text{ for each } z \in \text{GR}(p^b,m).$$
For $b=1,$ generalized trace map ${\Tr}$  reduces to the usual trace map  ${\tr}: \mathbb{F}_{p^m}  \to \mathbb{F}_{p}$  defined by ${\tr}(z) =z+z^p+z^{p^2}+\cdots+z^{p^{m-1}}$ for all $z \in \mathbb{F}_{p^m}.$ 
\begin{prop}\cite{sison}  The generalized trace map ${\Tr}:\text{GR}(p^b,m) \to  \mathbb{Z}_{p^b}$ has the following properties.
\begin{itemize}
\item ${\Tr}(u+v)={\Tr}(u)+{\Tr}(v)$ for each $u,v \in \text{GR}(p^b,m).$
\item  ${\Tr}(e u)=e{\Tr}(u)$ for each $u \in \text{GR}(p^b,m)$ and $e \in \mathbb{Z}_{p^b}.$
\item  ${\Tr}$ is surjective and $\text{GR}(p^b,m)/ \ker( {\Tr} )=\mathbb{Z}_{p^b}.$
%\item ${\Tr}(f(u))=f({\Tr}(u))={\Tr}(u)$ for each $u \in \text{GR}(p^b,m).$
\end{itemize}
\end{prop}


\begin{prop}\label{prop:charGR}\cite{shuqin} The map $\chi: \text{GR}(p^b,m) \to \mathbb{C}^*$ defined by
$\chi(r)=\zeta^{{\Tr}(r)}$, with $\zeta = \exp\bigl(\frac{2\pi i}{p^b}\bigr)$, is a generating character of  $\text{GR}(p^b,m).$
\end{prop}


\begin{defin} 
\label{def:trace_duals}
For an additive code $\mathcal{C}$ of $\text{GR}(p^b,m)^{2n},$  we define
\begin{itemize}
\item the trace-symplectic dual of $\mathcal{C}$ as
$$\mathcal{C}^{\perp_{{\Tr}}} = \{v \in \text{GR}(p^b,m)^{2n} : \Tr( \langle v | c\rangle_{\mathrm{s}}) = 0 \text{~for all~} c \in \mathcal{C}\},$$
\item and for $t = 0,1,\ldots,b$, the trace-symplectic $t$-dual of $\mathcal{C}$ as 
$$\mathcal{C}^{\perp_{\Tr,t}} = \{v \in \text{GR}(p^b,m)^{2n} ~:~ \Tr( \langle v | c\rangle_{\mathrm{s}}) \equiv 0 \!\!\! \pmod{p^{b-t}} \text{~for all~} c \in \mathcal{C}\}.$$
\end{itemize}
\end{defin}
For $b=1,$  $\mathcal{C}^{\perp_{{\Tr}}}$ is equal to 
$\mathcal{C}^{\perp_{{\tr}}} =\{v \in \mathbb{F}_{p^m}^{2n} ~:~ \tr(\langle v | c\rangle_{\mathrm{s}}) = 0 \text{ for all } c \in \mathcal{C}\}$.

\begin{prop}\label{prop:dualGR}  For an additive code $\mathcal{C}$ of $\text{GR}(p^b,m)^{2n},$ we have $\mathcal{C}^{\perp_\chi}=\mathcal{C}^{\perp_{{\Tr}}}$, and $\mathcal{C}^{\perp_{\chi,t}}=\mathcal{C}^{\perp_{\Tr,t}}$ for $t = 0,1,\ldots,b$.
\end{prop}
\begin{proof}Using  Proposition \ref{prop:charGR}, we observe that $\chi(r)=1$ if and only if ${\Tr}(r)=0,$ and $\chi(r) \in \langle \zeta^{p^{b-t}} \rangle $ if and only if ${\Tr}(r) \equiv 0 \pmod{p^{b-t}}$.  This implies that $\mathcal{C}^{\perp_\chi}=\mathcal{C}^{\perp_{{\Tr}}}$, and $\mathcal{C}^{\perp_{\chi,t}}=\mathcal{C}^{\perp_{\Tr,t}}.$
\end{proof} 

\begin{defin}\cite{sison} Two bases $\{\beta_1,\beta_2,\ldots,\beta_m\}$ and  $\{\gamma_1,\gamma_2,\ldots,\gamma_m\}$ of $\text{GR}(p^b,m)$ as a free module over $\mathbb{Z}_{p^b}$ are said to be dual if ${\Tr}(\beta_i\gamma_j)=\delta_{ij},$ where $\delta_{ij}$ denotes the Kronecker delta.\end{defin}
\begin{thm}\cite{sison} Every basis $\{\beta_1,\beta_2,\ldots,\beta_m\}$  of $\text{GR}(p^b,m)$ as a free module over $\mathbb{Z}_{p^b}$ has a unique dual.
\end{thm}
 Next we provide an upper bound on the size of a symplectic subset of $\text{GR}(p^b,m)^{2n}.$
\begin{prop}\label{prop:symGR}  If $\{a_{11},a_{12},a_{21},a_{22},\ldots,a_{e1},a_{e2}\} \subseteq \text{GR}(p^b,m)^{2n}$ is a symplectic subset of $\text{GR}(p^b,m)^{2n},$ then $e \leq nm.$ Moreover, for any given $z_1,z_2,\ldots,z_e \in \mathbb{Z}_{p^b}$ with $1 \leq e \leq nm,$ there exists a symplectic subset $\{a_{11},a_{12},a_{21},a_{22},$ $\ldots,a_{e1},a_{e2}\} \subseteq \text{GR}(p^b,m)^{2n}$ such that $ \chi(\langle a_{j1} |  a_{j2}\rangle_{\text{s}})= e^{i \frac{2\pi }{p^b} z_j  }$ for $1\leq j \leq e.$
\end{prop}
\begin{proof} To prove this result, we first note that $\chi(r)=e^{i \frac{2\pi}{p^b}  {\Tr}(r)}$ for each $r \in \text{GR}(p^b,m),$ which implies that $\chi(r)=1$ if and only if ${\Tr}(r)=0.$ Let $\gamma_1,\gamma_2,\ldots,\gamma_m$ be a  basis of $\text{GR}(p^b,m)$ as a free module over $\mathbb{Z}_{p^b}$ and $\beta_1,\beta_2,\ldots,\beta_m$ be the dual basis of $\gamma_1,\gamma_2,\ldots,\gamma_m.$   Now define a map 
$\phi: \text{GR}(p^b,m) \to \mathbb{Z}_{p^b}^m$ as follows:
$$r=r_0+r_1\gamma+\cdots+r_{m-1}\gamma_{m-1}  \stackrel{\phi}{\longmapsto}  (r_1,r_2,\ldots,r_m).$$ 
%
Extend this map to $\phi': \text{GR}(p^b,m) ^{2n} \to \mathbb{Z}_{p^b}^{2nm}$ as follows:
$$(d_1,d_2,\ldots,d_n,e_1,e_2,\ldots,e_n) \stackrel{\phi'}{\longmapsto} (\phi(d_1),\phi(d_2),\ldots,\phi(d_n),\phi(e_1),\phi(e_2),\ldots,\phi(e_n)).$$
It can be checked that $${\Tr}(\langle (d_1,d_2,\ldots,d_n,e_1,e_2,\ldots,e_n) \mid (d'_1,d'_2,\ldots,d'_n,e'_1,e'_2,\ldots,e'_n) \rangle_{\text{s}} )~~~~~~~~~~~~~~~~~~~~~~~~~~~~~~~~~~~~~~~~~~~~~~~~~~$$ $$=\langle (\phi(d_1),\phi(d_2,)\ldots,\phi(d_n),\phi(e_1),\phi(e_2),\ldots,\phi(e_n)) \mid (\phi(d'_1),\phi(d'_2,),\ldots,\phi(d'_n),\phi(e'_1),\phi(e'_2),\ldots,\phi(e'_n)) \rangle_{\text{s}} \, . $$
%
From this, we see that if $\{a_{11},a_{12},a_{21},a_{22},\ldots,a_{e1},a_{e2}\} \subseteq \text{GR}(p^b,m)^{2n}$ is a symplectic subset of $\text{GR}(p^b,m)^{2n},$ then $\{\phi'(a_{11}),\phi'(a_{12}),\phi'(a_{21}),\phi'(a_{22}),\ldots,\phi'(a_{e1}),\phi'(a_{e2})\} \subseteq \mathbb{Z}_{p^b}^{2mn}$ is a symplectic subset of $\mathbb{Z}_{p^b}^{2mn}$. Hence, by Corollary \ref{cor2}, we have $e \leq nm.$ 

To prove the next part,  let $t $ be a non-negative integer such that $tm < e \leq (t+1)m$. Then, for $j = k \ell + m$, with $0 \le k \le t-1$ and $1 \le \ell \le m$, set 
$$
a_{j1} = (\underbrace{ 0,0,\ldots,0}_{k},z_{km+\ell}\gamma_{\ell},\underbrace{ 0,0,\ldots,0}_{2n-k-1}) \ \text{ and } \ a_{j2} = (\underbrace{ 0,0,\ldots,0}_{n+k},\beta_{\ell},\underbrace{ 0,0,\ldots,0}_{n-k-1}), 
$$
and for $j = tm + \ell$, with $1 \le \ell \le e-tm$, set
$$
a_{j1} = (\underbrace{ 0,0,\ldots,0}_{t},z_{tm+\ell}\gamma_{\ell},\underbrace{ 0,0,\ldots,0}_{2n-t-1}) \ \text{ and } \ a_{j2} = (\underbrace{ 0,0,\ldots,0}_{n+t},\beta_{\ell},\underbrace{ 0,0,\ldots,0}_{n-t-1}).
$$ 
Then, $\{a_{11},a_{12},a_{21},a_{22},$ $\ldots,a_{e1},a_{e2}\}$ is a symplectic subset of $\text{GR}(p^b,m)^{2n}$ with the required property.
\end{proof}

We can now obtain an explicit formula for the minimum entanglement degree of an additive code $\mathcal{C} \subseteq \text{GR}(p^b,m)^{2n}.$ 

\begin{thm}\label{EAQZ} Any minimal $\chi$-self-orthogonal extension of an additive code $\mathcal{C} \subseteq \text{GR}(p^b,m)^{2n}$ has entanglement degree equal to $\left\lceil { \frac{1}{2m} \rank(\mathcal{C}/(\mathcal{C} \cap \mathcal{C}^{\perp_{{\Tr}}}))  }\right\rceil.$ 
\end{thm}
\begin{proof} 
Let $c := \frac12 \rank(\mathcal{C}/(\mathcal{C} \cap \mathcal{C}^{\perp_{{\chi}}}))  $, which, by Proposition~\ref{prop:dualGR}, equals 
$\frac{1}{2}  \rank(\mathcal{C}/(\mathcal{C} \cap \mathcal{C}^{\perp_{{\Tr}}}))  $. Let $c_{\min}$ be the entanglement degree of a minimal $\chi$-self-orthogonal extension, $\mathcal{C}'$, of $\mathcal{C}$. We want to show that $c_{\min} = \lceil \frac1m c\rceil$.

From Theorem~\ref{sym}, $\mathcal{C}$ has a generating set with $c$ hyperbolic pairs. Then, by Proposition~\ref{prop:ext}, there is a symplectic subset of $\mathbb{Z}_{p^a}^{2c_{\min}}$ of cardinality $2c$, so that by Proposition~\ref{prop:symGR}, we have $c \le mc_{\min}$. Thus, $c_{\min} \ge \lceil \frac1m c\rceil$.

For the inequality in the opposite direction, set $\nu := \lceil \frac1m c\rceil$, so that $c \le m\nu$. Let $(v_j,w_j)$ and $(x_j,y_j)$, $j = 1,2,\ldots,c$, be the hyperbolic pairs in the generating set of $\mathcal{C}$, and let $z_j = {\Tr}(\langle (v_j,w_j) \mid (x_j,y_j)\rangle_{\text{s}})$, $j=1,2,\ldots,c$. By Proposition~\ref{prop:symGR}, there exists a symplectic subset $\{a_{11},a_{12},a_{21},a_{22},$ $\ldots,a_{c1},a_{c2}\} \subseteq \text{GR}(p^b,m)^{2\nu}$ such that $ \chi(\langle a_{j1} |  a_{j2}\rangle_{\text{s}}) = e^{i \frac{2\pi }{p^b} z_j  }$ for $1\leq j \leq c$. Since $e^{i \frac{2\pi }{p^b} z_j} = \chi(\langle (v_j,w_j) \mid (x_j,y_j)\rangle_{\text{s}})$, we have a symplectic subset satisfying the condition of Proposition~\ref{prop:ext}. Hence, there exists a $\chi$-self-orthogonal extension $\mathcal{C}'$ with entanglement degree $\nu$. Therefore, $c_{\min} \le \nu = \lceil \frac1m c\rceil$, completing the proof of the theorem.
\end{proof}

Theorem~\ref{dt3} is an immediate consequence of Theorems~\ref{tEAC}, \ref{thm:eaqecc},  and \ref{EAQZ}, again using the fact (Proposition \ref{prop:dualGR}) that $\mathcal{C}^{\perp_\chi} = \mathcal{C}^{\perp_{{\Tr}}}$ and $\mathcal{C}^{\perp_{\chi,t}}=\mathcal{C}^{\perp_{\Tr,t}}$. The theorem is restated below with the numbers $\rho_t$ made explicit.

\begin{thm}[Restatement of Theorem~\ref{dt3}] Let $\mathcal{C} \subseteq \text{GR}(p^b,m)^{2n}$ be an additive code over the Galois ring $\text{GR}(p^b,m)$.   From $\mathcal{C}$, we can construct an  $((n,K,D;c))$ EAQECC over $\text{GR}(p^b,m),$ where the minimum number, $c$, of entanglement qudits needed for the construction is equal to  $\left\lceil \frac1{2m}\rank(\mathcal{C}/(\mathcal{C} \cap \mathcal{C}^{\perp_{{\Tr}}}))  \right\rceil$, the minimum distance is 
\begin{equation*}
D~=~\left\{\begin{array}{ll}
d_{\mathrm{s}}(\mathcal{C}^{\perp_{{\Tr}}}) & \text{~~if~~} \mathcal{C}^{\perp_{{\Tr}}} \subseteq \mathcal{C} \\
d_{\mathrm{s}}(\mathcal{C}^{\perp_{{\Tr}}}\setminus \mathcal{C} )  & \text{~~otherwise}\,.
\end{array}\right.
\end{equation*}
and the dimension $K$ is bounded as $p^{bm(n+c)}/ (|\mathcal{C}| \, {p^{\sum_{t=1}^{b-1} (b-t)\rho_t}}) \leq  K \leq p^{bm(n+c)}/|\mathcal{C}|$, with $\rho_t := \rank(\mathcal{C}/(\mathcal{C} \cap \mathcal{C}^{\perp_{\Tr,t-1}}))-\rank(\mathcal{C}/(\mathcal{C} \cap \mathcal{C}^{\perp_{\Tr,t}}))$. Additionally, if $\mathcal{C}$ is free, then $K = p^{bm(n+c)}/|\mathcal{C}|$.
\end{thm}

Corollary~\ref{cor:Fpm} follows from the observations that, for finite fields $\F_{p^m}$, the generalized trace map $\Tr$ reduces to the usual trace map $\tr: \F_{p^m} \to \F_p$, and $\mathcal{C}$ and $\mathcal{C} \cap \mathcal{C}^{\perp_{\tr}}$ are vector spaces over $\F_p$, so that, being free modules, their rank equals their dimension over $\F_p$. Moreover, $\dim_{\F_p}(\mathcal{C}) = \dim_{\F_p}(\mathcal{C} \cap \mathcal{C}^{\perp_{\tr}}) + \dim_{\F_p}(\mathcal{C} / (\mathcal{C} \cap \mathcal{C}^{\perp_{\tr}}))$.


%\begin{thm} Let $\mathcal{C} \subseteq \text{GR}(p^b,m)^{2n}$ be an additive code.   From $\mathcal{C}$, we can construct an  $((n,p^{bm(n+c)}/{|\mathcal{C}|} ,D;c))$ EAQECC over $\text{GR}(p^b,m),$ where the minimum number, $c$, of entanglement qudits needed for the construction is equal to  $ \left\lceil  \frac{1}{2m} (\rank(\mathcal{C})-\rank(\mathcal{C} \cap \mathcal{C}^{\perp_{{\Tr}}} ))\right\rceil$, and
%\begin{equation*}D~=~\left\{\begin{array}{ll}
%d_{\mathrm{s}}(\mathcal{C}^{\perp_{\Tr}}\setminus \mathcal{C} \cap\mathcal{C}^{\perp_{\Tr}})  & \text{~~if~~} \mathcal{C}\neq\mathcal{C}^{\perp_{\Tr}} ;\\
%d_{\mathrm{s}}(\mathcal{C}^{\perp_{\Tr}}) & \text{~~if~~}\mathcal{C}=\mathcal{C}^{\perp_{\Tr}} .\end{array}\right. \vspace{-1mm}\end{equation*}
%\end{thm}


\section{Conclusion and Future Work}\label{sec6}
In this paper, we proposed a framework to construct, from first principles, EAQECCs from additive codes over finite commutative local Frobenius rings. Given an additive code $\mathcal{C}$ over such a ring $\mathcal{R}$, we provide a means of constructing an EAQECC that uses a certain number of pre-shared entanglement qudits, a number that is determined by the algebraic structure of the code $\mathcal{C}$. This yields an upper bound on the minimum number of entanglement qudits required to construct an EAQECC from $\mathcal{C}$. We show that this bound is in fact tight for EAQECCs constructed from additive codes over $\mathbb{Z}_{p^a}$, but it can be improved when we start from codes over more general Galois rings. We derive an explicit expression for the minimum number of pre-shared entanglement qudits required to construct an EAQECC from an additive code over a Galois ring. This result significantly extends known results for EAQECCs constructed from additive codes over finite fields.

An interesting direction of future work would be to extend our formula for the minimum number of pre-shared entanglement qudits to EAQECCs over general finite commutative local Frobenius rings (beyond Galois rings). It would also be useful to find constructions of EAQECCs over Frobenius rings with good parameters, for example, codes that saturate the generalized quantum Singleton bound applicable to EAQECCs \cite{GHW21}. Finally, it would be of considerable interest to extend the EAQECC framework to non-local Frobenius rings.

\begin{thebibliography}{99}
\bibitem{ashikhmin} {A. Ashikhmin and E. Knill}, \newblock{``Nonbinary quantum stabilizer codes"}, \newblock{\it IEEE Trans. Inf. Theory}, 47   (2001), pp.\ 3065--3072.
%\bibitem{tittel} {H. Bechmann-Pasquinucci and W. Tittel},  \newblock{``Quantum cryptography using larger alphabets"},  \newblock{\it Phys. Rev. A}, 61 (2000), 062308. 
\bibitem{brun}  {T. Brun, I. Devetak and M. Hsieh}, \newblock{``Correcting quantum errors with entanglement"}, \newblock{\it Science}, 314  (2006), pp.\ 436--439.
%\bibitem{calder} { A. R. Calderbank, E. M. Rains, P. Shor, N. J. Sloane}, {``Quantum error correction via codes over GF(4)"}, \newblock{\it  IEEE Trans. Inf. Theory}, 44 (1998), pp. 1369-1387.
\bibitem{calderbank} { A. R. Calderbank and  P. Shor}, {``Good quantum error-correcting codes exist"}, \newblock{\it  Phys. Rev. A}, 54 (1996), pp.\ 1098--1105.
\bibitem{galindo} { C. Galindo, F.  Hernando, R. Matsumoto and D. Ruano}, \newblock{``Entanglement-assisted quantum error-correcting codes over arbitrary finite fields"}, \newblock{\it Quantum Inf. Process.}, 18 (2019), 116.
%\bibitem{guen} { K. Guenda, S. Jitman, T.A. Gulliver}, \newblock{``Constructions of good entanglement-assisted quantum error correcting codes"}, \newblock{\it  Des. Codes Cryptogr.}, 86 (2018), pp. 121--136.
\bibitem{gottesman} {D. Gottesman}, \newblock{\it Stabilizer Codes and Quantum Error Correction}, \newblock{Ph.D.\ dissertation}, CalTech, 1997.
\bibitem{DG04} {D. Gottesman}, CO639 course Problem Set 1 and solutions, \url{https://www2.perimeterinstitute.ca/personal/dgottesman/CO639-2004/}
%\bibitem{lai} {C.-Y. Lai, T. Brun},  \newblock{``Entanglement increases the error-correcting ability of quantum error-correcting codes"},  \newblock{\it Phys. Rev. A}, 88 (2013), 012320.
\bibitem{lee} {S. Lee and A. Klappenecker}, \newblock{``Entanglement-assisted quantum error correcting codes over nice rings"}, in \emph{Proc.\ 52nd Annual Allerton Conference on Communication, Control, and Computing}, 2014, pp.\ 1362--1367.
%\bibitem{gedik} {Z. Gedik, I. A. Silva, B. \c{C}akmak, G. Karpat,  E. L. G. Vidoto, D. O. Soares-Pinto, E. R. deAzevedo and F. F. Fanchini}, \newblock{``Computational speed-up with a single qudit"} \newblock{\it Sci. Rep.}, 5 (2015), 14671.
\bibitem{luerssen} {H. Gluesing-Luerssen and T. Pllaha},  \newblock{``On quantum stabilizer codes derived from local Frobenius rings"},  \newblock{\it  Finite Fields Appl.}, 58 (2019) pp.\ 145--173.
\bibitem{GHW21} M.\ Grassl, F.\ Huber, and A.\ Winter, ``Entropic proofs of Singleton bounds for quantum error-correcting codes,'' arXiv:2010.07902.
\bibitem{ketkar} {A. Ketkar, A. Klappenecker, S. Kumar, and P. Sarvepalli}, \newblock{``Nonbinary stabilizer codes over finite fields"}, \newblock{\it IEEE Trans. Inf. Theory}, 52   (2006), pp. 4892-4914.
\bibitem{luo} {L.\ Luo, Z.\ Ma, Z.\ Wei, and R.\ Leng}, \newblock{``Non-binary entanglement-assisted quantum stabilizer codes''}, \newblock{\it Sci.
China Inf. Sci.}, 60 (2017), 42501.
%\bibitem{mahmoudi} {S. Mahmoudi and K. Samei,}   \newblock{``Additive codes over Galois rings"},  \newblock{\it  Finite Fields Appl.}, 56 (2019), pp.\ 332--350.
\bibitem{nadella} {S. Nadella and A. Klappenecker}, \newblock{``Stabilizer codes over Frobenius rings"}, in \newblock{\it Proc.\ 2012 IEEE International Symposium on Information Theory (ISIT 2012)}, pp.\ 165--169.
\bibitem{nadkarni} {P. J. Nadkarni, S. S. Garani}, \newblock{``Non-binary entanglement-assisted stabilizer codes"}, \newblock{\it  Quantum Inf. Process.}, 20 (2021), 256. 
\bibitem{shor} {P. W. Shor},  \newblock{``Scheme for reducing decoherence in quantum computer memory"},  \newblock{\it Phys. Rev. A}, 52 (1995), pp.\ 2493--2496.
\bibitem{shuqin} {F. Shuqin and H. Wenbao},  \newblock{``Character sums over Galois rings and primitive polynomials over finite fields"},  \newblock{\it  Finite Fields Appl.}, 10 (2004), pp.\ 36--52.
\bibitem{sison} {V. P. Sison}, \newblock{``Bases and automorphism matrix of the Galois ring $GR(p^r, m)$ over $\mathbb{Z}_{p^r}"$}, \newblock{\it International Electronic Journal of Algebra}, 28 (2020), pp.\ 206--219.
\bibitem{steane} {A. Steane},  \newblock{``Error-correcting codes in quantum theory"},  \newblock{\it Phys. Rev. Lett.}, 77 (1996), pp.\ 793--797. 
\bibitem{wan} {Z.-X. Wan}, \newblock{\it Lectures on Finite Fields and Galois Rings}, \newblock{World Sci. Publ. Co. Inc.}, 2003.
\bibitem{wilde} {M. M. Wilde and T. A. Brun},  \newblock{``Optimal entanglement formulas for entanglement-assisted quantum coding"},  \newblock{\it Phys. Rev. A}, 77 (2008), 064302.
\end{thebibliography}


\end{document}
